\phantomsection
\addcontentsline{toc}{section}{Apuntes}
\begin{center}
\end{center}

\phantomsection
\addcontentsline{toc}{subsection}{Unidad 1. Fundamentos de Química}
\subsection*{\textcolor{blueunam}{Unidad 1. Fundamentos de Química}}

\phantomsection
\addcontentsline{toc}{subsubsection}{Clase 1 -- 18/08/2026: Definiciones y conceptos básicos}
\subsubsection*{\textcolor{blueunam}{Clase 1 -- 18/08/2026: Definiciones y conceptos básicos}}

Reactivo: Sustancia que puede interactuar o reaccionar con otra sustancia. (Se transforma --> Proceso de una sustancia para formar otra.) Un producto también puede actuar como reactivo A + B $\longleftrightarrow$  C $\rightarrow$ D \\
Sustrato o Materia prima: es la especie o reactivo específico sobre el cual actúa otro compuesto (un reactivo o catalizador) para transformarse en un producto, mientras que la materia prima química es el insumo base general. \\

Ejemplo:
\begin{center}
\schemestart
\chemfig{*6(-=-=-=)}
\arrow{->[\ce{HNO3}][\ce{H2SO4}]}
\chemfig{*6(-=(-NO_2)-=-=)}
\schemestop
\end{center}

El sustrato o Materia prima es la que se conserva en este caso la molécula de benceno es la que se conserva \\
Reactivo: Benceno, ácido nítrico y ácido sulfúrico \\
Producto: Nitrobenceno \\

Todo sustrato es un Reactivo, pero no todo Reactivo es un Sustrato. \\

\begin{center}
\schemestart
\chemname{\chemfig{*6(-=*5(-N(-H)-==-)-=-=-)}}{Indol}
\arrow{->[\shortstack{\ce{NBS}\\\small reactivo}][\shortstack{\ce{DMF}\\\small medio}]}
\chemname{\chemfig{*6(-=*5(-N(-H)-==-)-=(-Br)-=-)}}{5-bromoindol}
\+
\chemname{\chemfig{*6(-=*5(-N(-H)-=(-Br)=-)-=(-Br)-=-)}}{3,5-dibromoindol}
\schemestop
\end{center}

Producto: Sustancia final de la reacción ``(Es de interés)'' \\
Subproducto: Sustancia residual o no deseada de la reacción \\

\phantomsection
\addcontentsline{toc}{subsubsection}{Clase 2 -- 20/08/2026: Especies reactivas}
\subsubsection*{\textcolor{blueunam}{Clase 2 -- 20/08/2026: Especies reactivas}}

A $\rightarrow$ B (70\,\%) + C (20\,\%) + D (10\,\%) \\
Producto de interés $\neq$ Producto principal \\ 

Especies Reactivas: Especies que aparecen durante la reacción, pero desaparece al final. \\

\begin{itemize}
    \item \textcolor{blueunam}{\textbf{1- Carbocationes}}
    \item \textcolor{blueunam}{\textbf{2- Carbaniones}}
    \item \textcolor{blueunam}{\textbf{3- Radicales}}
    \item \textcolor{blueunam}{\textbf{4- Carbenos}}
\end{itemize}

\textcolor{blueunam}{\textbf{1- Carbocationes}}

\begin{center}
\schemestart
\chemfig{R-C^{\color{catred}+}(-[2]R)(-[6]R)}
\schemestop
\end{center}

\ce{e^{\color{anblue}-}} $-$ \ce{e^{\color{anblue}-}} libres $-$ \#enlaces $=$ Carga del átomo

Geometría (E.V.) \\

Repaso E.V.: 
Flujo de trabajo: \\
\begin{itemize}
    \item 1-Lewis
    \item 2-TRPECV
    \item 3-T.E.V.
    \item 4-T.O.M.
\end{itemize}

Lewis: \\
\begin{center}
\begin{tabular}{ll}
$\text{H}_2\text{O}$ & \\[6pt]
$\begin{array}{r}
2\,\ce{e^{\color{anblue}-}} \\
+\,6\,\ce{e^{\color{anblue}-}} \\
\hline
8\,\ce{e^{\color{anblue}-}} \\
-\,4\,\ce{e^{\color{anblue}-}} \\
\hline
4\,\ce{e^{\color{anblue}-}} \\
-\,4\,\ce{e^{\color{anblue}-}} \\
\hline
0\,\ce{e^{\color{anblue}-}}
\end{array}$
&
\chemfig{H-\charge{45=\:,135=\:}{O}-H}
\end{tabular}
\end{center}

\[
\begin{array}{rl}
\text{C.F. O} & = 6 - 2 - 4 = 0 \\
\text{C.F. H} & = (1 - 1 - 0) \times 2 = 0 \\
\hline
\text{Carga total} & = 0
\end{array}
\]

TRPECV: \\
Debido a que tiene 4 especies y 2 de ellas son pares de electrones
libres entonces la geometría es angular y podríamos sacar la hibridación
esperada que sería la $sp^3$ \\
\begin{center}
\chemfig{H-[104.5]\charge{45=\:,135=\:}{O}-[-104.5]H}
\end{center}

T.E.V.: \\
\ce{O}:$1s^2 2s^2 2p^4$ 


\begin{center}
\begin{tikzpicture}
\shorthandoff{>}

% --- Configuración original ---
% 1s
\draw (0,0) rectangle (0.8,0.8);
\node at (0.4,0.4) {$\upharpoonleft\downharpoonright$};
\node at (0.4,-0.3) {\small 1s};

% 2s
\draw (1.2,0) rectangle (2,0.8);
\node at (1.6,0.4) {$\upharpoonleft\downharpoonright$};
\node at (1.6,-0.3) {\small 2s};

% 2px
\draw (2.4,0) rectangle (3.2,0.8);
\node at (2.8,0.4) {$\upharpoonleft\downharpoonright$};
\node at (2.8,-0.3) {\small $2p_x$};

% 2py
\draw (3.4,0) rectangle (4.2,0.8);
\node at (3.8,0.4) {$\upharpoonleft$};
\node at (3.8,-0.3) {\small $2p_y$};

% 2pz
\draw (4.6,0) rectangle (5.4,0.8);
\node at (5,0.4) {$\upharpoonleft$};
\node at (5,-0.3) {\small $2p_z$};

% Flecha hibridación
\draw[-{Stealth}, thick] (5.8,0.4) -- (8.6,0.4);
\node at (7.2,0.7) {Hibridación};

% --- Configuración hibridada ---
% 1s (no cambia)
\draw (9.0,0) rectangle (9.8,0.8);
\node at (9.4,0.4) {$\upharpoonleft\downharpoonright$};
\node at (9.4,-0.3) {\small 1s};

% sp2 - caja 1
\draw (10.2,0) rectangle (11.0,0.8);
\node at (10.6,0.4) {$\upharpoonleft\downharpoonright$};
\node at (10.6,-0.3) {\small $sp^3$};

% sp2 - caja 2
\draw (11.4,0) rectangle (12.2,0.8);
\node at (11.8,0.4) {$\upharpoonleft\downharpoonright$};
\node at (11.8,-0.3) {\small $sp^3$};

% sp2 - caja 3
\draw (12.6,0) rectangle (13.4,0.8);
\node at (13.0,0.4) {$\upharpoonleft$};
\node at (13.0,-0.3) {\small $sp^3$};

% 2pz
\draw (13.8,0) rectangle (14.6,0.8);
\node at (14.2,0.4) {$\upharpoonleft$};
\node at (14.2,-0.3) {\small $sp^3$};

\end{tikzpicture}
\end{center}

Sé que es sp³ porque el átomo central tiene 4 grupos (2 pares libres + 2 enlaces). 
El número de grupos debe coincidir con el número de orbitales híbridos. 
El 1s no participa porque es un orbital de core (baja energía, no de valencia). 
Se hibridan 2s + 2px + 2py + 2pz: 1 orbital s y 3 orbitales p → sp³. \\

\textbf{La hibridación del carbono en el carbocatión:} \\

\begin{center}
\schemestart
\chemname{\chemfig{R-C^{\color{catred}+}(-[2]R)(-[6]R)}}{Carbocatión}
\schemestop
\end{center}

E.V. del carbocatión: \\
\ce{C}: $1s^2\,2s^2\,2p^2$ \textit{\textbf{(importante recordar que el carbono está cargado positivamente, así que la configuración cambia)}} $\rightarrow$ C$^{\color{catred}+}$: $1s^2\,2s^2\,2p^1$

\begin{center}
\begin{tikzpicture}
\shorthandoff{>}

% --- Fila 1: Original → Promoción ---
% 1s
\draw (0,0) rectangle (0.8,0.8);
\node at (0.4,0.4) {$\upharpoonleft\downharpoonright$};
\node at (0.4,-0.3) {\small 1s};

% 2s
\draw (1.2,0) rectangle (2,0.8);
\node at (1.6,0.4) {$\upharpoonleft\downharpoonright$};
\node at (1.6,-0.3) {\small 2s};

% 2px
\draw (2.4,0) rectangle (3.2,0.8);
\node at (2.8,0.4) {$\upharpoonleft$};
\node at (2.8,-0.3) {\small $2p_x$};

% 2py vacío
\draw (3.4,0) rectangle (4.2,0.8);
\node at (3.8,-0.3) {\small $2p_y$};

% 2pz vacío
\draw (4.6,0) rectangle (5.4,0.8);
\node at (5,-0.3) {\small $2p_z$};

% Flecha promoción
\draw[-{Stealth}, thick] (5.8,0.4) -- (8.6,0.4);
\node at (7.2,0.7) {Promoción};

% --- Configuración promovida ---
% 1s
\draw (9.0,0) rectangle (9.8,0.8);
\node at (9.4,0.4) {$\upharpoonleft\downharpoonright$};
\node at (9.4,-0.3) {\small 1s};

% 2s con 1 electrón
\draw (10.2,0) rectangle (11.0,0.8);
\node at (10.6,0.4) {$\upharpoonleft$};
\node at (10.6,-0.3) {\small 2s};

% 2px
\draw (11.4,0) rectangle (12.2,0.8);
\node at (11.8,0.4) {$\upharpoonleft$};
\node at (11.8,-0.3) {\small $2p_x$};

% 2py
\draw (12.6,0) rectangle (13.4,0.8);
\node at (13.0,0.4) {$\upharpoonleft$};
\node at (13.0,-0.3) {\small $2p_y$};

% 2pz vacío
\draw (13.8,0) rectangle (14.6,0.8);
\node at (14.2,-0.3) {\small $2p_z$};

% --- Fila 2: Promoción → Hibridación ---
% 1s
\draw (0,-2.5) rectangle (0.8,-1.7);
\node at (0.4,-2.1) {$\upharpoonleft\downharpoonright$};
\node at (0.4,-2.8) {\small 1s};

% 2s
\draw (1.2,-2.5) rectangle (2,-1.7);
\node at (1.6,-2.1) {$\upharpoonleft$};
\node at (1.6,-2.8) {\small 2s};

% 2px
\draw (2.4,-2.5) rectangle (3.2,-1.7);
\node at (2.8,-2.1) {$\upharpoonleft$};
\node at (2.8,-2.8) {\small $2p_x$};

% 2py
\draw (3.4,-2.5) rectangle (4.2,-1.7);
\node at (3.8,-2.1) {$\upharpoonleft$};
\node at (3.8,-2.8) {\small $2p_y$};

% 2pz vacío
\draw (4.6,-2.5) rectangle (5.4,-1.7);
\node at (5,-2.8) {\small $2p_z$};

% Flecha hibridación
\draw[-{Stealth}, thick] (5.8,-2.1) -- (8.6,-2.1);
\node at (7.2,-1.8) {Hibridación};

% --- Configuración hibridada sp2 ---
% 1s (no cambia)
\draw (9.0,-2.5) rectangle (9.8,-1.7);
\node at (9.4,-2.1) {$\upharpoonleft\downharpoonright$};
\node at (9.4,-2.8) {\small 1s};

% sp2 - caja 1
\draw (10.2,-2.5) rectangle (11.0,-1.7);
\node at (10.6,-2.1) {$\upharpoonleft$};
\node at (10.6,-2.8) {\small $sp^2$};

% sp2 - caja 2
\draw (11.4,-2.5) rectangle (12.2,-1.7);
\node at (11.8,-2.1) {$\upharpoonleft$};
\node at (11.8,-2.8) {\small $sp^2$};

% sp2 - caja 3
\draw (12.6,-2.5) rectangle (13.4,-1.7);
\node at (13.0,-2.1) {$\upharpoonleft$};
\node at (13.0,-2.8) {\small $sp^2$};

% 2pz vacío
\draw (13.8,-2.5) rectangle (14.6,-1.7);
\node at (14.2,-2.8) {\small $2p_z$};

\end{tikzpicture}
\end{center}

Geometría trigonal plana, 120° por ángulo \\


Electrófilo — especie que busca electrones (electron-lover).
Puede ser catión, molécula con carga parcial positiva o átomo
con orbital vacío. No necesariamente tiene carga formal positiva,
solo tiene deficiencia electrónica. \\

Nucleófilo — especie que dona electrones (nucleus-lover).
Tiene par de electrones disponible para compartir. Puede
ser anión, molécula con pares libres o enlace $\pi$. \\

\begin{center}
\begin{tikzpicture}
    % Tres enlaces a 120°
    \draw (0,0) -- (1.5,0) node[right] {R};
    \draw (0,0) -- (-0.75,1.3) node[above left] {R};
    \draw (0,0) -- (-0.75,-1.3) node[below left] {R};

    % Lóbulo superior
    \draw[dashed] (0,0.15) .. controls (-0.35,0.7) and (-0.35,1.3) .. (0,1.4)
                  .. controls (0.35,1.3) and (0.35,0.7) .. (0,0.15);

    % Lóbulo inferior
    \draw[dashed] (0,-0.15) .. controls (-0.35,-0.7) and (-0.35,-1.3) .. (0,-1.4)
                  .. controls (0.35,-1.3) and (0.35,-0.7) .. (0,-0.15);

    % Átomo central encima de todo
    \node[fill=white] at (0,0) {C$^{\color{catred}+}$};

    % Etiqueta
    \node[right] at (0.2,1.0) {\small $2p_z$ vacío};

\end{tikzpicture}

\small Carbocatión --- trigonal plana, $sp^2$, $2p_z$ vacío perpendicular al plano
\end{center}

\textbf{Efecto inductivo:} redistribución de densidad
electrónica a través de enlaces $\sigma$ por diferencia de
electronegatividad entre átomos. El átomo más electronegativo
polariza el enlace hacia sí, transmitiendo el efecto a lo largo
de la cadena con intensidad decreciente. Puede ser estabilizador
(grupos electrodonadores) o desestabilizador (grupos electroatractores)
según la dirección del jalón electrónico. \\

\textbf{Hiperconjugación:} deslocalización parcial de electrones
de un enlace $\sigma$ (C--H o C--C) hacia un orbital $p$ vacío
adyacente mediante solapamiento lateral. Requiere alineación paralela
entre el enlace $\sigma$ y el orbital $p$. Estabiliza al sistema al
aumentar la deslocalización electrónica, reduciendo la energía total.
No es un enlace $\pi$ completo sino una interacción $\sigma \rightarrow p$
de tres centros y dos electrones. \\

\begin{center}
\begin{tikzpicture}[scale=1.4]
\shorthandoff{>}

% --- Carbocatión C+ ---
\draw (0,0) -- (-0.75,1.3) node[above left] {H};
\draw (0,0) -- (-0.75,-1.3) node[below left] {H};

% Lóbulo superior pz
\draw[dashed] (0,0.15) .. controls (-0.35,0.7) and (-0.35,1.3) .. (0,1.4)
                  .. controls (0.35,1.3) and (0.35,0.7) .. (0,0.15);

% Lóbulo inferior pz
\draw[dashed] (0,-0.15) .. controls (-0.35,-0.7) and (-0.35,-1.3) .. (0,-1.4)
                  .. controls (0.35,-1.3) and (0.35,-0.7) .. (0,-0.15);

% Enlace C+--CH3
\draw (0,0) -- (2.0,0);
\node[fill=white] at (0,0) {C$^{\color{catred}+}$};

% --- CH3 adyacente ---
\draw (2.0,0) -- (2.0,1.3) node[above] {H};
\draw (2.0,0) -- (2.0,-1.3) node[below] {H};
\draw (2.0,0) -- (3.3,0) node[right] {H};
\node[fill=white] at (2.0,0) {C};

% Etiqueta pz
\node[right] at (0.2,1.0) {\small $2p_z$ vacío};

% Flecha 1: C-H superior → lóbulo superior pz
\draw[-{Stealth}, thick] (2.0,1.1) .. controls (1.2,1.8) and (0.4,1.6) .. (0.1,1.0);

% Flecha 2: C-H inferior → lóbulo inferior pz
\draw[-{Stealth}, thick] (2.0,-1.1) .. controls (1.2,-1.8) and (0.4,-1.6) .. (0.1,-1.0);

\node at (1.0,2.0) {\small Hiperconjugación};

% Puntos enlace C+-C (centrados en el enlace)
\filldraw (0.9,0) circle (1.2pt);
\filldraw (1.1,0) circle (1.2pt);

% Puntos enlace C-H superior (centrados en el enlace)
\filldraw (2.0,0.6) circle (1.2pt);
\filldraw (2.0,0.8) circle (1.2pt);

% Puntos enlace C-H inferior
\filldraw (2.0,-0.6) circle (1.2pt);
\filldraw (2.0,-0.8) circle (1.2pt);

% Puntos enlace C-H lateral
\filldraw (2.6,0) circle (1.2pt);
\filldraw (2.8,0) circle (1.2pt);

\end{tikzpicture}
\end{center}

Tipos de Carbocatión: \\

\begin{center}
\textbf{Carbocatión Primario} \\[6pt]
\begin{tikzpicture}
    % Tres enlaces a 120°
    \draw (0,0) -- (1.5,0) node[right] {H};
    \draw (0,0) -- (-0.75,1.3) node[above left] {$CH_3$};
    \draw (0,0) -- (-0.75,-1.3) node[below left] {H};

    % Lóbulo superior
    \draw[dashed] (0,0.15) .. controls (-0.35,0.7) and (-0.35,1.3) .. (0,1.4)
                  .. controls (0.35,1.3) and (0.35,0.7) .. (0,0.15);

    % Lóbulo inferior
    \draw[dashed] (0,-0.15) .. controls (-0.35,-0.7) and (-0.35,-1.3) .. (0,-1.4)
                  .. controls (0.35,-1.3) and (0.35,-0.7) .. (0,-0.15);

    % Átomo central encima de todo
    \node[fill=white] at (0,0) {C$^{\color{catred}+}$};

    % Etiqueta
    \node[right] at (0.2,1.0) {\small $2p_z$ vacío};

\end{tikzpicture}

\small Carbocatión de tipo primario, tiene 3 interacciones
de hiperconjugación posibles, una por cada enlace C-H del $CH_3$,
que contribuyen por rotación libre del grupo metilo.
\end{center}

\begin{center}
\textbf{Carbocatión Secundario} \\[6pt]
\begin{tikzpicture}
    % Tres enlaces a 120°
    \draw (0,0) -- (1.5,0) node[right] {$CH_3$};
    \draw (0,0) -- (-0.75,1.3) node[above left] {$CH_3$};
    \draw (0,0) -- (-0.75,-1.3) node[below left] {H};

    % Lóbulo superior
    \draw[dashed] (0,0.15) .. controls (-0.35,0.7) and (-0.35,1.3) .. (0,1.4)
                  .. controls (0.35,1.3) and (0.35,0.7) .. (0,0.15);

    % Lóbulo inferior
    \draw[dashed] (0,-0.15) .. controls (-0.35,-0.7) and (-0.35,-1.3) .. (0,-1.4)
                  .. controls (0.35,-1.3) and (0.35,-0.7) .. (0,-0.15);

    % Átomo central encima de todo
    \node[fill=white] at (0,0) {C$^{\color{catred}+}$};

    % Etiqueta
    \node[right] at (0.2,1.0) {\small $2p_z$ vacío};

\end{tikzpicture}

\small Tiene 6 interacciones de hiperconjugación posibles.
\end{center}

\begin{center}
\textbf{Carbocatión Terciario} \\[6pt]
\begin{tikzpicture}
    % Tres enlaces a 120°
    \draw (0,0) -- (1.5,0) node[right] {$CH_3$};
    \draw (0,0) -- (-0.75,1.3) node[above left] {$CH_3$};
    \draw (0,0) -- (-0.75,-1.3) node[below left] {$CH_3$};

    % Lóbulo superior
    \draw[dashed] (0,0.15) .. controls (-0.35,0.7) and (-0.35,1.3) .. (0,1.4)
                  .. controls (0.35,1.3) and (0.35,0.7) .. (0,0.15);

    % Lóbulo inferior
    \draw[dashed] (0,-0.15) .. controls (-0.35,-0.7) and (-0.35,-1.3) .. (0,-1.4)
                  .. controls (0.35,-1.3) and (0.35,-0.7) .. (0,-0.15);

    % Átomo central encima de todo
    \node[fill=white] at (0,0) {C$^{\color{catred}+}$};

    % Etiqueta
    \node[right] at (0.2,1.0) {\small $2p_z$ vacío};

\end{tikzpicture}

\small Tiene 9 interacciones de hiperconjugación posibles.
\end{center}

\begin{center}
\textbf{Carbocatión metílico} \\[6pt]
\begin{tikzpicture}
    % Tres enlaces a 120°
    \draw (0,0) -- (1.5,0) node[right] {$H$};
    \draw (0,0) -- (-0.75,1.3) node[above left] {$H$};
    \draw (0,0) -- (-0.75,-1.3) node[below left] {$H$};

    % Lóbulo superior
    \draw[dashed] (0,0.15) .. controls (-0.35,0.7) and (-0.35,1.3) .. (0,1.4)
                  .. controls (0.35,1.3) and (0.35,0.7) .. (0,0.15);

    % Lóbulo inferior
    \draw[dashed] (0,-0.15) .. controls (-0.35,-0.7) and (-0.35,-1.3) .. (0,-1.4)
                  .. controls (0.35,-1.3) and (0.35,-0.7) .. (0,-0.15);

    % Átomo central encima de todo
    \node[fill=white] at (0,0) {C$^{\color{catred}+}$};

    % Etiqueta
    \node[right] at (0.2,1.0) {\small $2p_z$ vacío};

\end{tikzpicture}

\small Tiene 0 interacciones de hiperconjugación posibles.
\end{center}

Como nota, para nombrarlos solo hay que contar cuántos carbonos
hay enlazados directamente con el carbono que nos interesa, en este caso carbocatión. \\

Otros tipos serían:

\begin{center}
\textbf{Carbocatión Primario} \\[6pt]
\schemestart
\chemname[1pt]{\chemfig{CH_3-CH_2-CH_2^{\color{catred}+}}}{Propilcarbocatión}
\schemestop
\\[2pt]
\small Tiene 2 interacciones de hiperconjugación posibles.
\end{center}

Entonces de saber que la hiperconjugación para este tipo de moléculas hace que sean más estables y entre más interacciones de hiperconjugación tengamos podemos decir que es más estable, por lo que el orden de estabilidad sería el siguiente: \\
\textit{Estabilidad} : $T > S > P > M$

\textbf{Nota:} La hiperconjugación también se observa en aromáticos. En el tolueno:

\begin{center}
\chemname{\chemfig{*6(=-=-(-CH_3)=-)}}{Tolueno}
\end{center}

Los enlaces C--H del \ce{CH3} se deslocalizan hacia el sistema $\pi$ del anillo, aumentando la densidad electrónica y haciendo al anillo más reactivo hacia electrófilos (sustitución electrofílica aromática). \\

\textbf{Resonancia}: fenómeno donde los electrones $\pi$ o pares libres no están localizados en un enlace específico sino deslocalizados sobre varios átomos. La molécula real no es ninguna de las estructuras resonantes sino el híbrido de todas ellas.

\textbf{Estabilidad}: la deslocalización reduce la energía del sistema porque los electrones ocupan un volumen mayor, lo que es energéticamente favorable. A más estructuras resonantes equivalentes → más deslocalización → más estabilidad.

\begin{center}
\schemestart
\chemname{\chemfig{=[1]-[-1]C^{\color{catred}+}(-[9]H)-[6]H}}{Catión alílico}
\qquad
\chemname{\chemfig{*6(=-=(-[1]CH_2^{\color{catred}+})-=-=-)}}{Catión bencílico}
\schemestop
\end{center}

Resonancia del Catión alílico:
\begin{center}
\schemestart
\chemfig{@{a1}=_[@{db}::30]-[@{sb}::-60]@{dnl}{\color{red}C^{\color{catred}+}}}
\arrow{<->}
\chemfig{@{c1}-[@{ab}::30]=_[@{pb}::-60]@{c2}{\color{red}C^{\color{catred}+}}}
\schemestop
\chemmove{
\draw[->](db)..controls +(-80:5mm) and +(-90:5mm)..(sb);}
\chemmove{
\draw[<-](ab)..controls +(-80:5mm) and +(-90:5mm)..(pb);}
\end{center}

Híbrido de resonancia del Catión alílico: \\
\begin{center}
$\left[\chemfig{-[1,,,,ddbond={-}]-[-1,,,,ddbond={-}]}\right]^{\color{catred}+}$
\end{center}

Resonancia del Catión Bencilico:

\begin{center}
\begin{tikzpicture}
\shorthandoff{>}

% Esquina superior izquierda - \textit{orto}1
\node at (-4,4) {\chemfig{*6(=-=(-[1]{\color{red}+})-=-=-)}};

% Esquina superior derecha - \textit{orto}2
\node at (4,4) {\chemfig{*6(=-{\color{red}+}-(=[1])-=-)}};

% Esquina inferior izquierda - CH2+
\node at (-4,-4) {\chemfig{*6(-=-(=[1])-{\color{red}+}-=)}};

% Esquina inferior derecha - para
\node at (4,-4) {\chemfig{*6({\color{red}+}-=-(=[1])-=-)}};

% Híbrido en el centro
\node at (0,0) {$\chemleft[\chemfig{**6(---(-[1,,,,ddbond={+}])---)}\chemright]^{\color{catred}+}$};

% Flechas horizontales
\draw[<->] (-2,4) -- (2,4);
\draw[<->] (-2,-4) -- (2,-4);

% Flechas verticales
\draw[<->] (-4,2) -- (-4,-2);
\draw[<->] (4,2) -- (4,-2);

\end{tikzpicture}
\end{center}

Fenilo:
\begin{center}
    \chemfig{*6(=-={\color{red}+}-=-)}
\end{center}

Vinílico:
\begin{center}
\chemfig{H-[:320]C(-[:220]H)=C^{\color{catred}+}-[:40]H}
\end{center}

Con base en los ejemplos anteriores, el orden de estabilidad de los carbocationes es: \\
\textit{Estabilidad}: $3^\circ \thickapprox B > 2^\circ \thickapprox A > 1^\circ > M \thickapprox V \thickapprox F$

\phantomsection
\addcontentsline{toc}{subsubsection}{Clase 3 -- 25/08/2026: Especies Reactivas}
\subsubsection*{\textcolor{blueunam}{Clase 3 -- 25/08/2026: Especies Reactivas}}

\textcolor{blueunam}{\textbf{2- Carbaniones}}

\begin{center}
    \schemestart
    \chemfig{R-C^{\color{anblue}-}(-[:90]R)(-[:270]R)}
    \schemestop
\end{center}

La hibridación del carbono en el carbanión: \\

E.V. del carbanión: \\
\ce{C}: $1s^2\,2s^2\,2p^2$ \textit{\textbf{(importante recordar que el carbono está cargado negativamente, así que la configuración cambia)}} $\rightarrow$ \ce{C^{\color{anblue}-}}: $1s^2\,2s^2\,2p^3$


\begin{center}
\begin{tikzpicture}
\shorthandoff{>}

% --- Fila 1: Original → Promoción ---
% 1s
\draw (0,0) rectangle (0.8,0.8);
\node at (0.4,0.4) {$\upharpoonleft\downharpoonright$};
\node at (0.4,-0.3) {\small 1s};

% 2s
\draw (1.2,0) rectangle (2,0.8);
\node at (1.6,0.4) {$\upharpoonleft\downharpoonright$};
\node at (1.6,-0.3) {\small 2s};

% 2px
\draw (2.4,0) rectangle (3.2,0.8);
\node at (2.8,0.4) {$\upharpoonleft$};
\node at (2.8,-0.3) {\small $2p_x$};

% 2py vacío
\draw (3.4,0) rectangle (4.2,0.8);
\node at (3.8,0.4) {$\upharpoonleft$};
\node at (3.8,-0.3) {\small $2p_y$};

% 2pz vacío
\draw (4.6,0) rectangle (5.4,0.8);
\node at (5,0.4) {$\upharpoonleft$};
\node at (5,-0.3) {\small $2p_z$};

% Flecha promoción
\draw[-{Stealth}, thick] (5.8,0.4) -- (8.6,0.4);
\node at (7.2,0.7) {Hibridación};

% --- Configuración promovida ---
% 1s
\draw (9.0,0) rectangle (9.8,0.8);
\node at (9.4,0.4) {$\upharpoonleft\downharpoonright$};
\node at (9.4,-0.3) {\small 1s};

% 2s con 1 electrón
\draw (10.2,0) rectangle (11.0,0.8);
\node at (10.6,0.4) {$\upharpoonleft\downharpoonright$};
\node at (10.6,-0.3) {\small $sp^3$};

% 2px
\draw (11.4,0) rectangle (12.2,0.8);
\node at (11.8,0.4) {$\upharpoonleft$};
\node at (11.8,-0.3) {\small $sp^3$};

% 2py
\draw (12.6,0) rectangle (13.4,0.8);
\node at (13.0,0.4) {$\upharpoonleft$};
\node at (13.0,-0.3) {\small $sp^3$};

% 2pz vacío
\draw (13.8,0) rectangle (14.6,0.8);
\node at (14.2,0.4) {$\upharpoonleft$};
\node at (14.2,-0.3) {\small $sp^3$};

\end{tikzpicture}
\end{center}



\begin{center}
\schemestart
\chemfig{
\charge{90=\:}{C}
(-[:330]R)
(<[:210]R)
(<:[:270]R)
}
\schemestop

\small Carbanión --- piramidal, $sp^3$, par libre en orbital $sp^3$ \\
Geometría Pirámide triangular, ángulo $<109.5^\circ$
\end{center}

\textbf{Efecto inductivo:} redistribución de densidad electrónica a 
través de enlaces $\sigma$ por diferencia de electronegatividad. 
En carbaniones, grupos electroatractores (\ce{-Cl}, \ce{-F}, \ce{-NO2}) 
estabilizan al retirar el exceso de carga negativa. Grupos electrodonadores 
(\ce{-CH3}, \ce{-OR}) desestabilizan al aumentarla. El efecto decae con la distancia. \\

\textbf{Hiperconjugación:} en carbaniones desestabiliza. Los enlaces $\sigma$ 
adyacentes donan densidad electrónica hacia el orbital $sp^3$ que ya contiene 
el par libre, incrementando el exceso de carga negativa en lugar de aliviarlo. 
(Véase \hyperref[sec:anexo:tom]{Anexo A.1.8} para el análisis completo de los 
mecanismos en competencia.) \\

Ejemplo: \\

\begin{center}
\schemestart
\chemname{\chemfig{\charge{90=\:}{C}
(-[:330]CH_3)
(<[:210]CH_3)
(<:[:270]CH_3)}}{Trimetilcarbanión}
\qquad\qquad
\chemname{\chemfig{\charge{90=\:}{C}
(-[:330]Cl)
(<[:210]Cl)
(<:[:270]Cl)}}{Triclorometil carbanión}
\schemestop
\\

\small El más estable de estos dos es el Triclorometil carbanión, porque los cloros alivian el exceso de carga del carbono.
\end{center}

El efecto de resonancia estabiliza: \\

Resonancia del Anión alílico:
\begin{center}
\schemestart
\chemfig{@{a1}=_[@{db}::30]-[@{sb}::-60]@{dnl}{\color{anblue}C^{\color{anblue}-}}}
\arrow{<->}
\chemfig{@{c1}-[@{ab}::30]=_[@{pb}::-60]@{c2}{\color{anblue}C^{\color{anblue}-}}}
\schemestop
\chemmove{
\draw[->](db)..controls +(-80:5mm) and +(-90:5mm)..(sb);}
\chemmove{
\draw[<-](ab)..controls +(-80:5mm) and +(-90:5mm)..(pb);}
\end{center}

Híbrido de resonancia del Anión alílico: \\
\begin{center}
$\left[\chemfig{-[1,,,,ddbond={-}]-[-1,,,,ddbond={-}]}\right]^{\color{anblue}-}$
\end{center}


Resonancia del Anión Bencilico:

\begin{center}
\begin{tikzpicture}
\shorthandoff{>}

% Esquina superior izquierda - \textit{orto}1
\node at (-4,4) {\chemfig{*6(=-=(-[1]{\color{anblue}-})-=-=-)}};

% Esquina superior derecha - \textit{orto}2
\node at (4,4) {\chemfig{*6(=-{\color{anblue}-}-(=[1])-=-)}};

% Esquina inferior izquierda - CH2+
\node at (-4,-4) {\chemfig{*6(-=-(=[1])-{\color{anblue}-}-=)}};

% Esquina inferior derecha - para
\node at (4,-4) {\chemfig{*6({\color{anblue}-}-=-(=[1])-=-)}};

% Híbrido en el centro
\node at (0,0) {$\chemleft[\chemfig{**6(---(-[1,,,,ddbond={+}])---)}\chemright]^{\color{anblue}-}$};

% Flechas horizontales
\draw[<->] (-2,4) -- (2,4);
\draw[<->] (-2,-4) -- (2,-4);

% Flechas verticales
\draw[<->] (-4,2) -- (-4,-2);
\draw[<->] (4,2) -- (4,-2);

\end{tikzpicture}
\end{center}

\textbf{Orden de estabilidad de carbaniones:}

La estabilidad depende de dos factores principales:
\begin{itemize}
    \item \textbf{Resonancia:} deslocaliza la carga negativa sobre varios átomos. Muy estabilizante. Bencilo y alílico quedan arriba por este efecto.
    \item \textbf{Efecto inductivo:} compite entre dos mecanismos (Coulomb y Heisenberg, véase \hyperref[sec:anexo:tom:heisenberg]{Anexo A.1.8}). En grupos alquilo el neto es desestabilizante.
\end{itemize}

Entre alquílicos simples el orden es inverso al de los carbocationes: más grupos alquilo donando carga, menos estable.

\textit{Estabilidad}: $B \thickapprox A > M > 1^\circ > 2^\circ > 3^\circ > V \thickapprox F$

\textcolor{blueunam}{\textbf{3- Radicales}}

\begin{center}
    \schemestart
    \chemfig{R-\charge{0=\.}{C}(-[:90]R)(-[:270]R)}
    \schemestop
\end{center}

La hibridación del carbono en los radicales: \\

E.V. del radical: \\
$1s^2 2s^2 2p^2$ \\
\textbf{Nota:} La carga formal se determina respecto al átomo neutro. 
El carbono neutro tiene 6 electrones. El radical tiene 6 electrones 
(carga 0), el carbanión tiene 7 (carga $-1$) y el carbocatión tiene 
5 (carga $+1$). El electrón desapareado del radical ya era propio del 
carbono neutro por Hund, no es un electrón ganado.


\begin{center}
\begin{tikzpicture}
\shorthandoff{>}

% --- Fila 1: Original → Promoción ---
% 1s
\draw (0,0) rectangle (0.8,0.8);
\node at (0.4,0.4) {$\upharpoonleft\downharpoonright$};
\node at (0.4,-0.3) {\small 1s};

% 2s
\draw (1.2,0) rectangle (2,0.8);
\node at (1.6,0.4) {$\upharpoonleft\downharpoonright$};
\node at (1.6,-0.3) {\small 2s};

% 2px
\draw (2.4,0) rectangle (3.2,0.8);
\node at (2.8,0.4) {$\upharpoonleft$};
\node at (2.8,-0.3) {\small $2p_x$};

% 2py vacío
\draw (3.4,0) rectangle (4.2,0.8);
\node at (3.8,0.4) {$\upharpoonleft$};
\node at (3.8,-0.3) {\small $2p_y$};

% 2pz vacío
\draw (4.6,0) rectangle (5.4,0.8);
\node at (5,0.4) {};
\node at (5,-0.3) {\small $2p_z$};

% Flecha promoción
\draw[-{Stealth}, thick] (5.8,0.4) -- (8.6,0.4);
\node at (7.2,0.7) {Promoción};

% --- Configuración promovida ---
% 1s
\draw (9.0,0) rectangle (9.8,0.8);
\node at (9.4,0.4) {$\upharpoonleft\downharpoonright$};
\node at (9.4,-0.3) {\small 1s};

% 2s con 1 electrón
\draw (10.2,0) rectangle (11.0,0.8);
\node at (10.6,0.4) {$\upharpoonleft$};
\node at (10.6,-0.3) {\small 2s};

% 2px
\draw (11.4,0) rectangle (12.2,0.8);
\node at (11.8,0.4) {$\upharpoonleft$};
\node at (11.8,-0.3) {\small $2p_x$};

% 2py
\draw (12.6,0) rectangle (13.4,0.8);
\node at (13.0,0.4) {$\upharpoonleft$};
\node at (13.0,-0.3) {\small $2p_y$};

% 2pz vacío
\draw (13.8,0) rectangle (14.6,0.8);
\node at (14.2,0.4) {$\upharpoonleft$};
\node at (14.2,-0.3) {\small $2p_z$};

% --- Fila 2: Promoción → Hibridación ---
% 1s
\draw (0,-2.5) rectangle (0.8,-1.7);
\node at (0.4,-2.1) {$\upharpoonleft\downharpoonright$};
\node at (0.4,-2.8) {\small 1s};

% 2s
\draw (1.2,-2.5) rectangle (2,-1.7);
\node at (1.6,-2.1) {$\upharpoonleft$};
\node at (1.6,-2.8) {\small 2s};

% 2px
\draw (2.4,-2.5) rectangle (3.2,-1.7);
\node at (2.8,-2.1) {$\upharpoonleft$};
\node at (2.8,-2.8) {\small $2p_x$};

% 2py
\draw (3.4,-2.5) rectangle (4.2,-1.7);
\node at (3.8,-2.1) {$\upharpoonleft$};
\node at (3.8,-2.8) {\small $2p_y$};

% 2pz vacío
\draw (4.6,-2.5) rectangle (5.4,-1.7);
\node at (5,-2.1) {$\upharpoonleft$};
\node at (5,-2.8) {\small $2p_z$};

% Flecha hibridación
\draw[-{Stealth}, thick] (5.8,-2.1) -- (8.6,-2.1);
\node at (7.2,-1.8) {Hibridación};

% --- Configuración hibridada sp2 ---
% 1s (no cambia)
\draw (9.0,-2.5) rectangle (9.8,-1.7);
\node at (9.4,-2.1) {$\upharpoonleft\downharpoonright$};
\node at (9.4,-2.8) {\small 1s};

% sp2 - caja 1
\draw (10.2,-2.5) rectangle (11.0,-1.7);
\node at (10.6,-2.1) {$\upharpoonleft$};
\node at (10.6,-2.8) {\small $sp^2$};

% sp2 - caja 2
\draw (11.4,-2.5) rectangle (12.2,-1.7);
\node at (11.8,-2.1) {$\upharpoonleft$};
\node at (11.8,-2.8) {\small $sp^2$};

% sp2 - caja 3
\draw (12.6,-2.5) rectangle (13.4,-1.7);
\node at (13.0,-2.1) {$\upharpoonleft$};
\node at (13.0,-2.8) {\small $sp^2$};

% 2pz vacío
\draw (13.8,-2.5) rectangle (14.6,-1.7);
\node at (14.2,-2.1) {$\upharpoonleft$};
\node at (14.2,-2.8) {\small $2p_z$};

\end{tikzpicture}
\end{center}

\begin{center}
\schemestart
\chemfig{
\charge{90=\.}{C}
(-[:330]R)
(<[:210]R)
(<:[:270]R)
}
\schemestop


\small Su geometría oscila entre Triangular y Tetraédrica y depende de los sustituyentes. \\
\end{center}

Ejemplos de radicales y comparación de su estabilidad:

\begin{center}
\begin{tikzpicture}
    % Tres enlaces a 120°
    \draw (0,0) -- (1.5,0) node[right] {$F$};
    \draw (0,0) -- (-0.75,1.3) node[above left] {$CH_3$};
    \draw (0,0) -- (-0.75,-1.3) node[below left] {$CH_3$};

    % Lóbulo superior
    \draw[dashed] (0,0.15) .. controls (-0.35,0.7) and (-0.35,1.3) .. (0,1.4)
                  .. controls (0.35,1.3) and (0.35,0.7) .. (0,0.15);

    % Lóbulo inferior
    \draw[dashed] (0,-0.15) .. controls (-0.35,-0.7) and (-0.35,-1.3) .. (0,-1.4)
                  .. controls (0.35,-1.3) and (0.35,-0.7) .. (0,-0.15);

    % Átomo central encima de todo
    \node[fill=white] at (0,0) {\charge{90[yshift=20pt]=\.}{C}};

    \node at (0,-2.0) {\small 2-fluoropropan-2-ilo};
\end{tikzpicture}

\qquad
\begin{tikzpicture}
    % Tres enlaces a 120°
    \draw (0,0) -- (1.5,0) node[right] {$F$};
    \draw (0,0) -- (-0.75,1.3) node[above left] {$CH_3$};
    \draw (0,0) -- (-0.75,-1.3) node[below left] {$F$};

    % Lóbulo superior
    \draw[dashed] (0,0.15) .. controls (-0.35,0.7) and (-0.35,1.3) .. (0,1.4)
                  .. controls (0.35,1.3) and (0.35,0.7) .. (0,0.15);

    % Lóbulo inferior
    \draw[dashed] (0,-0.15) .. controls (-0.35,-0.7) and (-0.35,-1.3) .. (0,-1.4)
                  .. controls (0.35,-1.3) and (0.35,-0.7) .. (0,-0.15);

    % Átomo central encima de todo
    \node[fill=white] at (0,0) {\charge{90[yshift=20pt]=\.}{C}};

    \node at (0,-2.0) {\small 1,1-difluoroetan-1-ilo};
\end{tikzpicture}
\qquad
\begin{tikzpicture}
    % Tres enlaces a 120°
    \draw (0,0) -- (1.5,0) node[right] {$F$};
    \draw (0,0) -- (-0.75,1.3) node[above left] {$F$};
    \draw (0,0) -- (-0.75,-1.3) node[below left] {$F$};

    % Lóbulo superior
    \draw[dashed] (0,0.15) .. controls (-0.35,0.7) and (-0.35,1.3) .. (0,1.4)
                  .. controls (0.35,1.3) and (0.35,0.7) .. (0,0.15);

    % Lóbulo inferior
    \draw[dashed] (0,-0.15) .. controls (-0.35,-0.7) and (-0.35,-1.3) .. (0,-1.4)
                  .. controls (0.35,-1.3) and (0.35,-0.7) .. (0,-0.15);

    % Átomo central encima de todo
    \node[fill=white] at (0,0) {\charge{90[yshift=20pt]=\.}{C}};

    \node at (0,-2.0) {\small trifluorometilo};
\end{tikzpicture}

\end{center}

La estabilidad de los radicales depende de la competencia entre efectos electrónicos que pueden
estabilizar o desestabilizar según el contexto y los sustituyentes presentes. Antes de describirlos,
una aclaración terminológica importante: el término \textbf{resonancia} se usa con tres significados
distintos en química orgánica.
\begin{enumerate}
    \item \textit{Estructuras resonantes}: el doble enlace se desplaza entre posiciones (catión
    alílico, benceno). Es el concepto clásico de resonancia.
    \item \textit{Efecto de resonancia}: donación de un par libre hacia un orbital $p$ semiocupado
    o vacío por solapamiento lateral. No hay desplazamiento de doble enlace.
    \item \textit{Resonancia en \textsc{tom}}: todos los efectos de deslocalización electrónica
    son manifestaciones del mismo fenómeno fundamental, el solapamiento de orbitales. La distinción
    entre hiperconjugación, efecto de resonancia e inductivo es didáctica, no física.
\end{enumerate}

En la Teoría de Orbitales Moleculares (\textsc{tom}), todos estos efectos son el mismo fenómeno:
solapamiento de orbitales que extiende la nube electrónica del electrón desapareado hacia mayor
volumen. Por el principio de incertidumbre de Heisenberg, mayor volumen implica menor energía
cinética promedio del electrón, lo que estabiliza al sistema (véase
\hyperref[sec:anexo:tom]{Anexo~TOM}). El \textsc{somo}\footnote{\textsc{somo} (del inglés
\textit{Singly Occupied Molecular Orbital}): orbital molecular con exactamente un electrón. Distinto
del \textsc{homo} (lleno, dos electrones) y del \textsc{lumo} (vacío). Es el orbital del electrón
desapareado del radical. Energéticamente se ubica entre ambos. El estado semilleno es un estado
estable: añadirle un electrón lo acercaría a un carbanión (menos estable) y quitarle el único que
tiene lo convertiría en un carbocatión. Cualquier desviación del estado semilleno desestabiliza.} es
el orbital del electrón desapareado. \\

\textbf{Hiperconjugación positiva} ($\sigma \rightarrow$ \textsc{somo})\textbf{:} cuando la
geometría permite la alineación de los enlaces $\sigma$ adyacentes (C--H o C--C) con el
\textsc{somo}, el electrón desapareado se deslocaliza hacia esos enlaces, extendiendo su nube
electrónica hacia mayor volumen y bajando la energía cinética del sistema por Heisenberg. Los
enlaces C--C actúan como puentes; los electrones con efecto real son los de los enlaces C--H en
posición $\boldsymbol{\beta}$ al carbono radical (carbono directamente unido al carbono
radical).\footnote{Las posiciones $\alpha$ y $\beta$ son nomenclatura general en química orgánica, no
exclusiva de radicales. Posición $\alpha$: átomos directamente unidos al grupo funcional o centro
de interés (aquí, el carbono que porta el electrón desapareado). Posición $\beta$: átomos unidos
al carbono $\alpha$, es decir a dos enlaces del grupo funcional. En hiperconjugación los que
importan son los \textbf{hidrógenos $\beta$}: los H unidos al carbono $\beta$, que son los que se
alinean geométricamente con el \textsc{somo} y solapan con él. Resultados espectroscópicos y
teóricos indican que aproximadamente el 15\,\% del electrón desapareado en el radical etilo se
deslocaliza hacia los tres hidrógenos $\beta$.} \textbf{La geometría es prerequisito, no consecuencia:} sin alineación
no hay solapamiento. A mayor número de enlaces C--H en posición $\beta$ disponibles, mayor
deslocalización y mayor estabilidad.

\textbf{Nota sobre carbocationes:} en carbocationes el \textsc{somo} está vacío, por lo que la
hiperconjugación positiva estabiliza por dos mecanismos simultáneos: deslocalización de la nube
(Heisenberg) y dispersión de la carga positiva. En radicales el \textsc{somo} ya contiene un
electrón, por lo que el único mecanismo estabilizador es la deslocalización. Por eso las diferencias
de estabilidad T/S/P/M son menores en radicales que en carbocationes. 
(Véase \hyperref[sec:anexo:tom:somo]{Anexo A.1.9} y \hyperref[sec:anexo:tom:heisenberg]{Anexo A.1.8}.)\\

\textbf{Hiperconjugación negativa} (\textsc{somo} $\rightarrow \sigma^*$)\textbf{:} el electrón
del \textsc{somo} puede extenderse hacia orbitales $\sigma^*$ antienlazantes vacíos de los enlaces
C--H vecinos alineados. La dirección es opuesta a la positiva, pero el mecanismo estabilizador es
el mismo: el electrón ocupa mayor volumen y baja su energía cinética. El $\sigma^*$ está disponible como aceptor porque el $\sigma$ ya está lleno con 
2 electrones. Por el principio de exclusión de Pauli, no caben más electrones en el 
$\sigma$ enlazante, así que cualquier electrón adicional que quiera deslocalizarse 
hacia ese enlace debe ocupar el $\sigma^*$ antienlazante. La hiperconjugación negativa contribuye a la
geometría ligeramente piramidal de los radicales alquilo simples.\footnote{Investigación reciente
(\textit{J. Org. Chem.}, 2026) demuestra mediante cálculos NBO que para radicales alquilo simples,
la hiperconjugación negativa es principalmente responsable de su geometría ligeramente piramidal.
La diferencia de energía entre la forma plana y piramidal es pequeña, por lo que en la práctica se
tratan como $sp^2$.} \\

\textbf{Efecto inductivo:} un átomo electronegativo en posición $\alpha$ (directamente unido al
carbono radical) atrae la densidad electrónica del \textsc{somo} hacia sí por diferencia de
electronegatividad, polarizando la nube sin cambio de orbital. A diferencia de la deslocalización,
la polarización \textit{confina} al electrón en un volumen menor, aumentando su energía cinética
promedio por Heisenberg y desestabilizando al radical. El electrón no se transfiere al átomo
electronegativo: sigue perteneciendo al carbono radical, pero su densidad de probabilidad se
desplaza hacia él.\footnote{Distinción clave: deslocalización = el electrón cambia de orbital y
ocupa mayor volumen, energía cinética baja, estabiliza. Polarización = el electrón permanece en el
mismo orbital pero su distribución espacial se contrae hacia el átomo electronegativo, energía
cinética sube, desestabiliza en radicales. En carbaniones el inductivo sí estabiliza porque alivia
el exceso de carga negativa, mecanismo distinto.} \\

\textbf{Donación por resonancia del par libre} (efecto de resonancia)\textbf{:} cualquier átomo
con pares libres puede donar al \textsc{somo} por solapamiento lateral, extendiendo la nube del
electrón desapareado hacia mayor volumen. El mecanismo estabilizador es idéntico a la
hiperconjugación: deslocalización, no donación de carga neta. El \textsc{somo} semilleno no
necesita ni se beneficia de carga adicional; lo que lo estabiliza es que su electrón ocupe mayor
volumen. Requiere alineación geométrica del par libre con el orbital $p$ del \textsc{somo}. \\

\textbf{Caso específico: F en posición $\beta$}\footnote{Este caso es distinto a los diagramas
mostrados arriba, donde el F está en posición $\alpha$ (directamente unido al carbono radical).
Aquí se describe una molécula tipo \ce{R-\dot{C}(-R)-CH2F}, donde el F está en el carbono
$\beta$.} Cuando el F está en posición $\beta$, coexisten dos efectos sobre el radical: la
donación del par libre del F al \textsc{somo} por resonancia (estabiliza, extiende la nube) y el
efecto inductivo del F (desestabiliza, contrae la nube). El resultado neto depende de cuál predomina.
Con uno o dos F en $\beta$, la resonancia domina. A medida que aumentan los F, el inductivo
acumulado puede superar a la resonancia. La geometría piramidal que genera el inductivo reduce
además el solapamiento del par libre con el \textsc{somo}, retroalimentando la desestabilización. \\

\textbf{Caso de los diagramas: F en posición $\alpha$} En los tres radicales de los diagramas
anteriores, el F está directamente unido al carbono radical (posición $\alpha$). En esta
configuración:
\begin{itemize}
    \item El par libre del F no puede donar efectivamente al \textsc{somo} en posición
    $\alpha$ porque el enlace C--F es $\sigma$ y el par libre del F está en el plano
    de los enlaces $\sigma$, mientras que el \textsc{somo} es un orbital $p$
    perpendicular a ese plano. La geometría no permite el solapamiento lateral
    necesario para la donación por resonancia.
    \item El efecto inductivo del F desestabiliza al radical por confinamiento de la nube.
    \item La hiperconjugación de los grupos CH$_3$ en posición $\beta$ estabiliza.
    \item El resultado neto depende de cuántos F hay vs cuántos H$_\beta$ disponibles.
\end{itemize}

Cuando coexisten efectos en la misma especie, el resultado neto depende de cuál predomina. Regla
práctica del curso: comparar cuántas interacciones de cada tipo están disponibles. El efecto con
más interacciones domina. El análisis en casos mixtos es una pregunta abierta que depende de la
intensidad relativa de cada efecto. \\

\textbf{Geometría y carácter $s$:} la geometría del radical oscila entre plana y piramidal según
los sustituyentes. El electrón desapareado repele menos que un par completo, por lo que el radical
tiende más a plano que el carbanión, pero no completamente. Cuando la deslocalización domina, el
electrón se extiende y su efecto de repulsión sobre los enlaces no está concentrado, favoreciendo
geometría plana (mayor carácter $s$, tendencia $sp^2$). Cuando el inductivo domina, el electrón
se contrae y su repulsión es más localizada, favoreciendo geometría piramidal (menor carácter $s$,
tendencia $sp^3$). Hablar de $sp^2$ o $sp^3$ para radicales es una aproximación didáctica: la
geometría real es un punto en ese continuo determinado por el balance de efectos electrónicos. \\

Para radicales hidrocarbonados (sin halógenos), donde solo opera la hiperconjugación:

\textit{Estabilidad}: $T > S > P > M$

\phantomsection
\addcontentsline{toc}{subsubsection}{Clase 4 -- 27/08/2026: Tipos de reacciones orgánicas}
\subsubsection*{\textcolor{blueunam}{Clase 4 -- 27/08/2026: Tipos de reacciones orgánicas}}

Sustituciones sobre C $sp^3$: \\
Tipos de sustituciones: \\

\begin{center}
\begin{tabular}{|l|l|l|l|}
\hline
\textbf{Tipo} & \textbf{Mecanismo} & \textbf{Subtipos} & \textbf{Cinética} \\
\hline
Iónica & Heterolítica & $S_N2$ & Concertada, bimolecular \\
\cline{3-4}
       &               & $S_N1$ & Por pasos, unimolecular \\
\hline
Radicalaria & Homolítica & Iniciación & Unimolecular \\
\cline{3-4}
            &            & Propagación & Bimolecular \\
\cline{3-4}
            &            & Terminación & Bimolecular \\
\hline
Pericíclica & Concertada cíclica & --- & Unimolecular \\
\hline
\end{tabular}
\end{center}

\textbf{Bimolecular:} la velocidad depende de la concentración de dos 
especies. \textbf{Unimolecular:} depende solo de una. \textbf{Concertada:} 
un solo paso, sin intermediario entre reactivos y productos. \textbf{Por 
pasos:} hay al menos un intermediario real entre reactivos y productos.

\textbf{Intermediarios.} En $S_N1$ el intermediario es el carbocatión. En 
reacciones radicalarias es el radical, generado en la iniciación y consumido 
en la terminación. En $S_N2$ y en pericíclicas no hay intermediario, solo un 
estado de transición. Tanto los carbocationes como los radicales ya se 
describieron con detalle en las clases anteriores.

\textbf{Ruptura heterolítica:} los dos electrones del enlace C--X se van con el halógeno,
formando X$^{\color{anblue}-}$. El carbono queda con deficiencia electrónica. Es el tipo de ruptura que ocurre en $S_N1$ y $S_N2$. \\
\textbf{Ruptura homolítica:} cada átomo se lleva un electrón del enlace, formando dos radicales.
Es la ruptura característica de las reacciones radicalarias. \\

Haluros de alquilo como sustratos: \\

\begin{center}
\schemestart
\chemname{\chemfig{CH_3-Br}}{Metílico}
\qquad
\chemname{\chemfig{CH_3-CH_2-Br}}{Primario}
\qquad
\chemname{\chemfig{CH_3-CH(-[2]CH_3)-Br}}{Secundario}
\qquad
\chemname{\chemfig{CH_3-C(-[2]CH_3)(-[6]CH_3)-Br}}{Terciario}
\schemestop
\end{center}

\textcolor{blueunam}{\textbf{Sustitución Nucleofílica Bimolecular ($S_N2$)}} \\

Mecanismo de un solo paso. El nucleófilo ataca por el lado opuesto al halógeno
(\textit{ataque anti}) simultáneamente mientras el halógeno se va como X$^{\color{anblue}-}$. Nunca
se forma un carbocatión libre. \\

% --- Línea 1: Reactivos → Estado de transición ---
\begin{center}
\schemestart
$Na^{\color{catred}+}$
\+
$\underbrace{\textcolor{anblue}{OH^{\color{anblue}-}}}_{\text{Nucleófilo}}$
\quad
\chemname{
    \chemfig{
        \overset{\delta^{\color{catred}+}}{C}(-[2]H)(-[6]H)(-[180]H)-[0]\underset{\delta^{\color{anblue}-}}{Br}
    }
}{Sustrato}
\arrow{->}
$\underbrace{
    \left[\,
        \textcolor{anblue}{HO}
        \cdots
        \chemfig{C(-[90]H)(-[270]H)(-[45]H)}
        \cdots
        Br
    \,\right]^{\ddagger}
}_{\text{Estado de transición}\footnotemark}$
\schemestop
\end{center}

% --- Línea 2: Estado de transición → Productos ---
\begin{center}
\schemestart
$\left[\,
    \textcolor{anblue}{HO}
    \cdots C \cdots
    Br
\,\right]^{\ddagger}$
\arrow{->}
\chemname{
    \chemfig{
        \textcolor{anblue}{HO}-[:0]C(-[:90]H)(<[:270]H)(<:[:150]H)
    }
}{Inversión\footnotemark}
\+
$\textcolor{anblue}{Br^{\color{anblue}-}}$
\+
$Na^{\color{catred}+}$
\schemestop
\end{center}

\footnotetext[1]{\textbf{Estado de transición ($\ddagger$):} especie hipotética de máxima
energía que no puede ser aislada. Representa el instante en que el enlace C--O se está
formando y el C--Br se está rompiendo simultáneamente. Los puntos ($\cdots$) indican
enlaces parciales. Es el máximo de la coordenada de reacción.}

\footnotetext[2]{\textbf{Inversión de Walden:} el nucleófilo ataca por el lado opuesto
al grupo saliente. Los tres sustituyentes del carbono se invierten como un paraguas que
se voltea. Si el carbono es quiral, el producto tiene configuración R/S opuesta al reactivo.}

\textbf{Quiralidad:} un carbono es quiral cuando tiene cuatro sustituyentes diferentes.
Solo en ese caso la inversión de configuración produce un estereoisómero distinto.
(Véase \hyperref[sec:anexo:quiralidad:quiral]{Anexo A.2.2} para la discusión completa.) \\

\textbf{Favorecida en:} haluros metílicos y primarios (menor impedimento estérico). \\

\textbf{Impedimento estérico:} obstáculo al acceso de un reactivo a una 
posición o átomo de la molécula generado por la presencia de átomos, grupos 
o nubes electrónicas que bloquean la aproximación. Tiene dos componentes 
que suelen ir juntos:
\begin{itemize}
\item \textbf{Geométrico:} el grupo ocupa espacio físico impidiendo que el 
reactivo se aproxime lo suficiente para que ocurra el solapamiento orbital 
necesario para la reacción.
\item \textbf{Electrónico:} grupos grandes y polarizables generan nubes 
electrónicas difusas que repelen al reactivo antes de que llegue al sitio 
de reacción.
\end{itemize}
A mayor tamaño del sustituyente, mayor impedimento estérico y menor 
reactividad en esa posición.



\begin{center}
\includegraphics[width=0.75\textwidth]{/home/azquepchup/IQM/Quimica_de_Materiales/Imagenes_QM/coordenada_SN2.pdf}
\end{center}

\small La gráfica muestra la coordenada de reacción para una sustitución exotérmica
(línea continua) y endotérmica (línea punteada). El símbolo $\ddagger$ indica el estado
de transición. $E_a$ es la energía de activación y $\Delta H$ indica si la reacción
libera o absorbe energía neta. \\

\normalsize
\textcolor{blueunam}{\textbf{Sustitución Nucleofílica Unimolecular ($S_N1$)}} \\

Mecanismo de dos pasos. En el primer paso (lento, determinante) el halógeno se va solo
formando un carbocatión como intermedio real. En el segundo paso (rápido) el nucleófilo
ataca al carbocatión. \\

Ejemplos de intermedios carbocatión: \\

\begin{center}
\schemestart
\chemname{\chemfig{CH_3-C^{\color{catred}+}(-[2]H)(-[6]H)}}{\footnotesize Metilcarbocatión}
\qquad
\qquad
\chemname{\chemfig{CH_3-C^{\color{catred}+}(-[2]CH_3)(-[6]H)}}{\footnotesize Isopropilcarbocatión}
\qquad
\qquad
\chemname{\chemfig{CH_3-C^{\color{catred}+}(-[2]CH_3)(-[6]CH_3)}}{\footnotesize Terc-butilcarbocatión}
\schemestop
\end{center}

\textbf{Favorecida en:} haluros terciarios (el carbocatión terciario es más estable
por mayor hiperconjugación). No hay inversión de configuración definida: el carbocatión
es plano ($sp^2$) y el nucleófilo puede atacar por cualquier lado. \\

\textbf{Nucleofilía} — capacidad de una especie para donar un par de electrones a un
centro electrofílico: \\

\begin{center}
\begin{tabular}{cccccccc}
$\underbrace{OH^{\color{anblue}-} \quad HS^{\color{anblue}-} \quad NH_2^{\color{anblue}-}}_{\text{más fuertes (carga negativa)}}$
& $H_2O$ & $R_3P$ & $R$--$OH$ & $NH_3$ & $H_2S$
\end{tabular}
\end{center}

Todos son nucleofílicos por carga negativa o pares libres. Los más fuertes tienen carga
negativa concentrada. Entre los cargados negativamente, el más fuerte es el que tiene la
carga más polarizable (menos electronegativo y más grande). \\

\textbf{Electrófilo:} especie con deficiencia de electrones. \\
\textbf{Nucleófilo:} especie con par de electrones disponible para donar. \\
\textbf{Sitio electrofílico:} zona de mayor densidad de carga positiva en la molécula,
donde el nucleófilo ataca.

\phantomsection
\addcontentsline{toc}{subsubsection}{Clase 5 -- 01/09/2026: Sustituciones en C $sp^3$}
\subsubsection*{\textcolor{blueunam}{Clase 5 -- 01/09/2026: Sustituciones en C $sp^3$}}

\begin{itemize}
\item Sustitución Nucleofílica Bimolecular ($S_N2$):
    haluros de alquilo primarios, secundarios y metílicos.
\item Sustitución Nucleofílica Unimolecular ($S_N1$):
    haluros de alquilo terciarios y secundarios, bencílicos y alílicos.
\end{itemize}

$S_N2$: No se forma carbocatión porque el carbono primario o metílico
no tiene suficiente estabilización para sostener esa carga,
prefieren pasar por el estado de transición, se llama Bimolecular
porque en el estado de transición participan dos moléculas, el nucleófilo y el sustrato
y sustitución porque se sustituye un grupo por otro. \\


$S_N1$: Se forma un carbocatión, es Unimolecular porque el paso determinante es la formación del carbocatión, y sustitución porque se sustituye un grupo por otro. \\

\begin{center}
\schemestart
\chemname{\chemfig{CH_3-C(-[:90]CH_3)(-[:0]x)(-[6]CH_3)}}{\footnotesize Haluro de alquilo terciario}
\qquad
\qquad
\chemname{\chemfig{=[:30]-[:-30]-[:30]x}}{\footnotesize Haluro alílico primario}
\qquad
\qquad
\chemname{\chemfig{*6(=-=(-[:30](-[:-30]x))-=-)}}{\footnotesize Haluro bencílico primario}
\schemestop
\end{center}

\begin{center}
\schemestart
$Na^{\color{catred}+}$
\+
$\underbrace{\textcolor{anblue}{OH^{\color{anblue}-}}}_{\text{Nucleófilo}}$
\quad
\chemname{
    \chemfig{
        \overset{\delta^{\color{catred}+}}{C}(-[2]H)(-[6]H)(-[180]H)-[0]\underset{\delta^{\color{anblue}-}}{x}
    }
}{\footnotesize Sustrato \\ \scriptsize (Ruptura heterolítica)}
\arrow{->}
$\underbrace{
    \chemfig{C^{\color{catred}+}(-[:90]H)(-[:270]H)(-[:180]H)} + \textcolor{anblue}{x^{\color{anblue}-}}
}_{\text{Intermediario}}$
\schemestop
\end{center} 



\begin{center}
\schemestart
\chemfig{C(-[:0]OH)(-[:90]H)(-[:270]H)(-[:180]H)}
\+ $Na^{\color{catred}+}$
\+ $\textcolor{anblue}{x^{\color{anblue}-}}$
\schemestop
\end{center}

\includegraphics[width=0.8\textwidth]{/home/azquepchup/IQM/Quimica_de_Materiales/Imagenes_QM/coordenada_SN1.pdf}

$S_N1$ genera mezclas racémicas porque el carbocatión intermedio es plano ($sp^2$):
el nucleófilo puede atacar con igual probabilidad por la cara superior o inferior
del orbital $p$ vacío, produciendo ambos enantiómeros R y S en proporción aproximadamente
1:1.\footnotemark
\footnotetext{Véase \hyperref[sec:anexo:quiralidad:rac]{Anexo A.2.4} para mezclas racémicas.}

\phantomsection
\addcontentsline{toc}{subsubsection}{Clase 6 -- 03/09/2026: Reacciones $S_N1$ y $S_N2$}
\subsubsection*{\textcolor{blueunam}{Clase 6 -- 03/09/2026: Reacciones $S_N1$ y $S_N2$}}

\begin{center}
\small
\renewcommand{\arraystretch}{1.6}
\setlength{\LTleft}{0pt}
\begin{longtable}{|p{2.5cm}|p{5cm}|p{5cm}|}
\hline
\rule{0pt}{14pt} & \textbf{$S_N1$} & \textbf{$S_N2$} \\
\hline
\endhead
\hline
\textbf{Sustrato} &
Terciario, bencílico, alílico. El carbocatión formado debe ser estable. &
Metílico, primario, secundario. Carbono accesible para el ataque del nucleófilo. \\
\hline
\textbf{Disolvente} &
Polar prótico (agua, alcoholes). Solvatiza el carbocatión y el grupo saliente,
favoreciendo la ionización y separando el par iónico. &
Polar aprótico (acetona, DMF, DMSO). No solvatiza al nucleófilo, lo deja libre y reactivo. \\
\hline
\textbf{Nucleófilo} &
Débil o fuerte (no importa). El paso lento es la ionización, no el ataque. &
Fuerte y poco voluminoso. Debe atacar al carbono central en un solo paso. \\
\hline
\textbf{Estereo-\linebreak química} &
Mezcla racémica $(\pm)$. El carbocatión plano ($sp^2$) puede ser atacado por ambas caras.
Con par iónico íntimo puede haber exceso enantiomérico parcial. &
Inversión de configuración (inversión de Walden). El ataque anti produce el enantiómero opuesto. \\
\hline
\textbf{Mecanismo} &
Dos pasos: (1)~ionización lenta $\rightarrow$ carbocatión (paso determinante);
(2)~ataque del nucleófilo rápido.
Unimolecular: velocidad $= k[\text{sustrato}]$. &
Un solo paso: el nucleófilo ataca mientras el grupo saliente se va simultáneamente.
Bimolecular: velocidad $= k[\text{sustrato}][\text{Nu}]$. \\
\hline
\textbf{Estados de\linebreak transición} &
Dos TS: TS$_1$ (ionización, paso lento) y TS$_2$ (ataque nucleofílico, paso rápido).
El carbocatión es el intermediario real entre ambos. &
Un único TS pentacoordinado (complejo activado de Walden): nucleófilo entra y grupo saliente
sale simultáneamente. No hay intermediario. \\
\hline
\end{longtable}
\normalsize
\end{center}


$S_N1$:
\begin{center}
\schemestart
\chemname{\chemfig{(-[2])(-[6])(-[4])-I}}{\footnotesize Haluro terciario}
\+
\ce{NaCN}
\arrow{->[\footnotesize \chemfig{C^{\color{catred}+}(-[2])(-[6])(-[4])} \+ $\textcolor{anblue}{I^{\color{anblue}-}}$][\footnotesize \ce{H2O}]}
\chemname{\chemfig{(-[2])(-[6])(-[4])-CN}}{\footnotesize Producto}
\+
\ce{NaI}
\schemestop
\end{center}


\begin{center}
\small
\renewcommand{\arraystretch}{1.4}
\begin{tabular}{|p{4cm}|p{4cm}|p{4cm}|}
\hline
\textbf{Polar prótico} & \textbf{Polar aprótico} & \textbf{Aprótico no polar} \\
\hline
Agua (\ce{H2O}) & Acetona (\ce{Me2CO}) & Hexano \\
Metanol (MeOH) & DMF (dimetilformamida) & Benceno \\
Etanol (EtOH) & DMSO (dimetilsulfóxido) & Tolueno \\
Isopropanol (iPrOH) & Acetonitrilo (MeCN) & Éter dietílico (\ce{Et2O}) \\
Ácido acético (AcOH) & THF (tetrahidrofurano) & \\
Ácido fórmico (HCOOH) & Diclorometano (DCM) & \\
\hline
\multicolumn{1}{|p{4cm}|}{\small Donan H, forman puentes de hidrógeno.
Favorecen \textbf{$S_N1$}.} &
\multicolumn{1}{p{4cm}|}{\small Sin H donables. No solvatan al nucleófilo.
Favorecen \textbf{$S_N2$}.} &
\multicolumn{1}{p{4cm}|}{\small Baja polaridad. No favorecen reacciones iónicas.} \\
\hline
\end{tabular}
\normalsize
\end{center}

\textbf{Solvatación}

La solvatación es el proceso por el cual las moléculas del disolvente rodean 
a una especie iónica o polar, estabilizándola electrostáticamente. Para un 
catión, el disolvente orienta su parte negativa (dipolo $\delta^{\color{anblue}-}$) hacia él; 
para un anión, orienta su parte positiva ($\delta^{\color{catred}+}$).

En disolventes \textbf{polares próticos} (agua, alcoholes), la solvatación 
es especialmente efectiva porque el H unido a O o N tiene carga parcial 
positiva $\delta^{\color{catred}+}$ que forma \textbf{puentes de hidrógeno} directamente 
con los aniones:

\[
\text{X}^{\color{anblue}-} \cdots \text{H--O--R}
\]

El H es el átomo más pequeño que existe: puede acercarse íntimamente al 
anión sin impedimento estérico, rodeándolo y estabilizándolo fuertemente. 
Esto no puede hacerlo cualquier $\delta^{\color{catred}+}$: el $\delta^{\color{catred}+}$ del C en la 
acetona o del S en el DMSO está en átomos más grandes que no pueden 
acercarse al anión con esa intimidad. Por eso solo el prótico estabiliza 
eficazmente los aniones.

Esta interacción específica y fuerte envuelve al anión con una capa de 
moléculas de disolvente, reduciendo su reactividad como nucleófilo. 
En disolventes \textbf{polares apróticos} (acetona, DMF, DMSO), no hay 
H disponible para el puente de hidrógeno: la solvatación del anión es 
débil, el nucleófilo queda ``desnudo'' y mucho más reactivo.

\textbf{Resumen:}
\begin{itemize}
\item \textbf{Prótico:} estabiliza aniones por puente de H (H desnudo 
$\delta^{\color{catred}+}$) y cationes por $\delta^{\color{anblue}-}$ del O. Nucleófilo solvatado = 
menos reactivo. Favorece $S_N1$.
\item \textbf{Aprótico polar:} estabiliza cationes por $\delta^{\color{anblue}-}$ pero 
no solvatiza bien los aniones. Nucleófilo libre = más reactivo. 
Favorece $S_N2$.
\end{itemize}

\textbf{Nota:} ``prótico'' no se refiere a que el protón haga algo activo. 
Se refiere a que el H desnudo por electronegatividad genera un $\delta^{\color{catred}+}$ 
lo suficientemente pequeño y accesible para formar puentes de hidrógeno 
con aniones. Es electrostática pura mediada por el tamaño del H. \\

\begin{center}
\includegraphics[width=0.85\textwidth]{/home/azquepchup/IQM/Quimica_de_Materiales/Imagenes_QM/solvatar.png}
\end{center}

\textbf{Pares iónicos en $S_N1$: íntimo, separado y libre}

En $S_N1$, cuando el grupo saliente X$^{\color{anblue}-}$ abandona el sustrato no siempre se aleja
completamente de inmediato. Dependiendo de la velocidad de solvatación del disolvente,
se pueden formar tres situaciones distintas en secuencia:

\begin{enumerate}
\item \textbf{Par iónico íntimo} (contacto): el X$^{\color{anblue}-}$ permanece muy cerca del carbocatión,
    bloqueando parcialmente una de sus caras. El nucleófilo ataca preferentemente por el
    lado opuesto, produciendo un exceso enantiomérico (resultado más parecido a $S_N2$ 
    que a una racémica perfecta).

\item \textbf{Par iónico separado} (solvent-separated): el disolvente se ha intercalado
    entre el carbocatión y el X$^{\color{anblue}-}$, separándolos parcialmente. La selectividad facial
    disminuye pero no desaparece. Se produce mezcla con algo de exceso enantiomérico.

\item \textbf{Carbocatión libre}: el X$^{\color{anblue}-}$ se ha alejado completamente. El carbocatión
    es plano ($sp^2$) y simétrico. El nucleófilo ataca con igual probabilidad por ambas
    caras: \textbf{mezcla racémica perfecta} (50\% R, 50\% S).
\end{enumerate}

\begin{center}
\small
\begin{tabular}{|p{3cm}|p{3.5cm}|p{3.5cm}|}
\hline
\textbf{Especie} & \textbf{Descripción} & \textbf{Resultado estereoquímico} \\
\hline
Par iónico íntimo & X$^{\color{anblue}-}$ junto al C$^{\color{catred}+}$, sin disolvente entre ellos &
Exceso enantiomérico parcial \\
\hline
Par iónico separado & Disolvente intercalado entre X$^{\color{anblue}-}$ y C$^{\color{catred}+}$ &
Mezcla con algo de selectividad \\
\hline
Carbocatión libre & X$^{\color{anblue}-}$ completamente alejado & Mezcla racémica $(\pm)$ 50:50 \\
\hline
\end{tabular}
\normalsize
\end{center}

El disolvente polar prótico acelera la separación del par iónico porque solvatiza
rápidamente tanto al carbocatión como al X$^{\color{anblue}-}$. Cuanto más rápida es la separación,
más se acerca el resultado a la racémica perfecta. Por eso en disolventes próticos
muy ionizantes como el agua, $S_N1$ produce mezclas casi perfectamente racémicas.

Grupo saliente: Un grupo saliente (abreviado como $LG$, del inglés \textit{leaving group}) es un átomo
o fragmento molecular que se separa de la molécula original durante una reacción de sustitución
($S_N1$ o $S_N2$) o de eliminación ($E1$ o $E2$), llevándose consigo el par de electrones del enlace
$\sigma$ que lo unía al carbono.

La calidad de un grupo saliente está determinada por su capacidad para estabilizar la carga negativa
que adquiere al salir. Cuanto más débil sea su base conjugada, mejor será como grupo saliente.
Esto se cuantifica mediante el $pK_a$ del ácido conjugado ($H-LG$): a menor $pK_a$, mejor es el grupo
saliente. Por ejemplo, el yoduro ($I^{\color{anblue}-}$) es un excelente grupo saliente porque su ácido conjugado ($HI$)
tiene un $pK_a \approx -10$, mientras que el hidróxido ($OH^{\color{anblue}-}$) es un pésimo grupo saliente porque su ácido
conjugado ($H_2O$) tiene un $pK_a \approx 15.7$.

Además de la basicidad, la polarizabilidad y la estabilización por resonancia facilitan la salida, tal como se explica a continuación:

\begin{center}
    \schemestart
    \chemname{\chemfig{(-[:0]x)(-[:90])(-[:180])(-[:270])}}{Haluro de alquilo}
    \schemestop
\end{center}
Los grupos salientes son mejores si son muy polarizables, esto se puede saber con la diferencia de tamaño
entre más tamaño más polarizable, por ejemplo el yodo es mejor que el bromo y este mejor que el cloro, la 
polarizabilidad estabiliza sus electrones lo que hace que tenga menos energía de enlace y más fácil de romper
por lo que son mejores grupos salientes\\

También se pueden tener grupos salientes con carga negativa como por ejemplo:
\begin{center}
    \schemestart
    \chemname{\chemfig{O^{\color{anblue}-}-S(=[:90]O)(=[:270]O)(-([:0]*6(=-=-=-)))}}{Tosilato}
    \schemestop
\end{center}
Estos se estabilizan con resonancia lo que lo hace más estable
y son mejores grupos salientes que los haluros, 
por ejemplo el grupo tosilato es mejor que el bromo.\\

Cinética:
$V = k[R-x][Nu]$
$k$: Constante cinética
$R-x$:Haluro de alquilo
$Nu$:Nucleofilo

Para la reacción $S_N2$ la velocidad depende 
de la concentración del nucleofilo y del haluro de alquilo, 
mientras que para la reacción $S_N1$ la velocidad depende 
solo de la concentración del haluro de alquilo.\\

Esto se debe a que el mecanismo de la reacción $S_N2$ es concertado:
el ataque del nucleófilo y la salida del grupo saliente ocurren en 
un solo paso a través de un estado de transición. En ese estado de 
transición, tanto el haluro de alquilo como el nucleófilo están involucrados 
simultáneamente, por lo que la velocidad depende de la concentración de ambos 
reactivos (es una reacción bimolecular).

En cambio, en la reacción $S_N1$ el mecanismo ocurre en dos pasos. El paso lento 
(etapa determinante de la velocidad) es la ruptura del enlace $C-X$ para formar 
un carbocatión y el grupo saliente. En este paso solo participa el haluro de alquilo; 
el nucleófilo aún no ha intervenido. El paso rápido es el ataque del nucleófilo al carbocatión. 
Como el paso lento solo involucra al haluro de alquilo, la concentración del nucleófilo no afecta 
la velocidad global, por lo que su ecuación cinética es $V = k[R-X]$ (reacción unimolecular). \\

\textbf{Para una discusión más profunda sobre polarizabilidad, nucleofilicidad 
y la base orbital de $S_N2$, véase} (Véase \hyperref[sec:anexo:nucleofilicidad:polar]{Anexo A.3.1} para polarizabilidad, 
\hyperref[sec:anexo:nucleofilicidad:carga]{Anexo A.3.2} para nucleofilicidad, y 
\hyperref[sec:anexo:nucleofilicidad:homo]{Anexo A.3.5} para la base orbital de SN2.).

\phantomsection
\addcontentsline{toc}{subsubsection}{Clase 7 -- 08/09/2026: Sustitución Electrofílica Aromática}
\subsubsection*{\textcolor{blueunam}{Clase 7 -- 08/09/2026: Sustitución Electrofílica Aromática}}

Los compuestos aromáticos son notablemente poco reactivos frente a reactivos electrofílicos en 
comparación con sus análogos no aromáticos. Un ejemplo ilustrativo es la reacción con \ce{Br2}:
el benceno no reacciona con \ce{Br2} bajo condiciones normales, mientras que su análogo lineal 
con el mismo número de carbonos e hidrógenos, el 1,3,5-hexatrieno, sí reacciona por adición 
electrofílica directa.

\begin{center}
\schemestart
\chemfig{*6(=-=-=-)}
\+
\ce{Br2}
\arrow{->}
\text{No reacciona}
\schemestop
\end{center}
\begin{center}
\schemestart
\chemfig{=[:30]-[:-30]=[:30]-[:-30]=[:30]}
\+
\ce{Br2}
\arrow{->}
\chemfig{(-[:270]Br)-[:30](-[:90]Br)-[:-30](-[:270]Br)-[:30](-[:90]Br)-[:-30](-[:270]Br)-[:30](-[:90]Br)}
\schemestop
\end{center}

Para que una molécula sea aromática tiene que cumplir con 4 criterios:
\begin{enumerate}
\item \textbf{Cíclica:} los átomos deben formar un anillo cerrado.
\item \textbf{Conjugación:} cada átomo del anillo debe aportar un orbital $p$ 
perpendicular al plano, ya sea por doble enlace, par libre (pirrol) u orbital 
vacío (cationes aromáticos). El patrón enlace--no enlace es la forma más común 
pero no la única.
\item \textbf{Plana:} todos los átomos del anillo deben estar en el mismo plano 
para que el solapamiento lateral de los orbitales $p$ sea efectivo.
\item \textbf{Regla de Hückel:} el número de electrones $\pi$ debe satisfacer 
$4n + 2$ con $n \in \mathbb{Z}_{\geq 0}$:
\[
n = \frac{e^{\color{anblue}-}_\pi - 2}{4}
\]
\end{enumerate}

\begin{center}
\renewcommand{\arraystretch}{2.0}
\begin{tabular}{|C{3.5cm}|C{1.8cm}|C{2cm}|C{1.8cm}|C{1.8cm}|C{2cm}|}
\hline
\textbf{Molécula} & \textbf{Cíclica} & \textbf{Conjugada} & \textbf{Plana} & \textbf{Hückel} & \textbf{Aromática} \\
\hline
\shortstack{1,3-butadieno \\ \small\chemfig{=[:30]-[:-30]=[:30]}} & No & Sí & No & No & No \\
\hline
\shortstack{Ciclopentadieno \\ \small\chemfig{*5(-=-=-)}} & Sí & No & No & No & No \\
\hline
\shortstack{Benceno \\ \small\chemfig{*6(=-=-=-)}} & Sí & Sí & Sí & Sí & Sí \\
\hline
\end{tabular}
\end{center}


\textbf{Notas:}

\textbf{Observación sobre $n$:} en policíclicos aromáticos fusionados simples 
(naftaleno, antraceno), $n$ en la regla de Hückel coincide numéricamente con 
el número de anillos fusionados. No es una definición general sino una 
coincidencia de estos casos.

\textbf{Planaridad geométrica vs planaridad electrónica:} son dos conceptos distintos.
\begin{itemize}
\item \textbf{Geométrica:} todos los átomos en el mismo plano en un instante dado. 
Un enlace $\sigma$ puede rotar y sacar átomos de ese plano.
\item \textbf{Electrónica:} todos los átomos del anillo tienen orbitales $p$ 
perpendiculares al plano disponibles para el sistema $\pi$. Depende de la 
hibridación, no de la rotación instantánea.
\end{itemize}
El criterio de planaridad en aromaticidad se refiere a la \textbf{electrónica}: 
los orbitales $p$ deben estar alineados para solapar. En práctica esto exige 
también planaridad geométrica, por eso ambas van juntas en aromáticos.

\textbf{Rotación del enlace $\sigma$:} es un proceso dinámico. La molécula no adopta 
una conformación fija sino una mezcla estadística ponderada por energía. Las 
conformaciones más estables se populan más. A temperatura baja se pueden ``congelar'' 
y detectar por separado; a temperatura ambiente interconvierten tan rápido que solo 
se observa el promedio. \\ 

\textbf{Grupos activantes (donadores) y Grupos desactivantes (electroatractores)} \\

DG = Grupo donador
WG = Grupo electroatractor

\begin{center}
\chemfig{*6(=-=-(-DG)=-)} \\
\quad
DG = 
\chemfig{\charge{90=\:,270=\:}{-O}H},
\chemfig{\charge{90=\:}{-N}H_2},
\chemfig{\charge{90=\:,270=\:}{-O}R},
\chemfig{\charge{90=\:}{-N}R_2},
\chemfig{\charge{90=\:,270=\:}{-S}H},
\chemfig{\charge{90=\:,270=\:}{-O}CH_3},
\chemfig{-CH_3}
\end{center}

\begin{center}
\chemfig{*6(=-=-(-WG)=-)} \\
\quad
WG = \ce{-CN}, \chemfig{-C(=[:90]O)-H}, \chemfig{-C(=[:90]O)-R}, 
\chemfig{-N^{\textcolor{catred}{+}}(=[:90]O)-O^{\textcolor{anblue}{-}}}
\end{center}

\textbf{Activaciones y desactivaciones y donde se realizan} \\

\begin{center}
\chemfig{*6(C_5=C_4-C_3=C_2-C_1=C_6-)}
\end{center}
\begin{center}
\renewcommand{\arraystretch}{1.4}
\begin{tabular}{|C{2cm}|C{2cm}|C{4cm}|}
\hline
\textbf{Posición} & \textbf{Nombre} & \textbf{Descripción} \\
\hline
1 & \textit{ipso} & Carbono que porta el sustituyente \\
\hline
2 y 6 & \textit{orto} & Adyacentes al $C_1$ \\
\hline
3 y 5 & \textit{meta} & Separados por un C del $C_1$ \\
\hline
4 & \textit{para} & Opuesto al $C_1$ \\
\hline
\end{tabular}
\end{center}
\textbf{Nota:} la numeración indica la posición del \textbf{carbono} en el anillo, 
no del sustituyente. El $C_1$ es el carbono que porta el sustituyente. La numeración 
parte del $C_1$ (carbono que porta el sustituyente) y avanza en cualquier dirección 
de forma creciente alrededor del anillo hasta regresar al $C_1$. Se elige la dirección
que dé los números más bajos a los demás sustituyentes presentes. \\

\begin{center}
\renewcommand{\arraystretch}{1.6}
\begin{tabular}{|L{3cm}|L{4cm}|L{4cm}|L{4cm}|}
\hline
\textbf{Propiedad} & \textbf{DG} & \textbf{WG} & \textbf{Halógenos} \\
\hline
\textbf{Mecanismo} & 
Donación de par libre al sistema $\pi$ del anillo por resonancia o hiperconjugación & 
Retiro de densidad $\pi$ del anillo por resonancia retractora o inductivo & 
Desactivan por inductivo (electronegatividad), donan por resonancia de par libre \\
\hline
\textbf{Efecto en anillo} & 
Activa, anillo más rico en electrones & 
Desactiva, anillo más pobre en electrones & 
Desactiva globalmente \\
\hline
\textbf{Dirige a} & 
\textit{orto} y \textit{para} & 
\textit{meta} & 
\textit{orto} y \textit{para} \\
\hline
\textbf{Ejemplos} & 
\ce{-OH}, \ce{-NH2}, \ce{-OR}, \ce{-NR2}, \ce{-CH3} & 
\ce{-NO2}, \ce{-CN}, \ce{-CHO}, \ce{-COR}, \ce{-COOH}, \ce{-CF3} & 
\ce{-F}, \ce{-Cl}, \ce{-Br}, \ce{-I} \\
\hline
\end{tabular}
\end{center}

\quad


\textbf{Ejemplo de un grupo activante:} \\
\schemestart
\chemname{\chemfig{*6(=-=-(-\charge{0=\:,180=\:}{O}-[:40]H)=-)}}{Fenol}
\arrow{<->}
\chemname{\chemfig{*6(=-=-(=\charge{180=\:}{\textcolor{catred}{O^{\color{catred}+}}}-[:40]H)-{\color{anblue}-}-)}}{\textit{orto}}
\arrow{<->}
\chemname{\chemfig{*6(={\color{anblue}-}-=-(=\charge{180=\:}{\textcolor{catred}{O^{\color{catred}+}}}-[:40]H)=-)}}{\textit{para}}
\arrow{<->}
\chemname{\chemfig{*6(=-={\color{anblue}-}-(=\charge{180=\:}{\textcolor{catred}{O^{\color{catred}+}}}-[:40]H)--)}}{\textit{orto}}
\schemestop \\

La densidad electrónica donada por el par libre del \ce{O} se deslocaliza 
por resonancia a través del sistema $\pi$ del anillo, enriqueciendo 
específicamente las posiciones \textit{orto} y \textit{para} por el patrón alternado 
de la conjugación. \\
\quad

\schemestart
\chemname{\chemfig{*6(=-=-(-N(=[:140]\charge{40=\:,140=\:}{O})-[:40]O)=-)}}{Nitrobenceno}
\arrow{<->}
\chemname{\chemfig{*6(=-=-(=N(-[:140]\charge{40=\:,140=\:,220=\:}{O^{\color{anblue}-}})-[:40]O)-{\color{catred}+}-)}}{Orto}
\arrow{<->}
\chemname{\chemfig{*6(-{\color{catred}+}-=-(=N(-[:140]\charge{40=\:,140=\:,220=\:}{O^{\color{anblue}-}})-[:40]O)-=)}}{Para}
\arrow{<->}
\chemname{\chemfig{*6(-=-{\color{catred}+}-(=N(-[:140]\charge{40=\:,140=\:,220=\:}{O^{\color{anblue}-}})-[:40]O)-=)}}{Orto}
\schemestop \\


Los WG no activan la posición \textit{meta}. La desactivan menos 
que las posiciones \textit{orto} y \textit{para}. El electrófilo prefiere el \textit{meta} no porque esté 
enriquecido sino porque es la posición relativamente menos empobrecida. \\

\phantomsection
\addcontentsline{toc}{subsubsection}{Clase 8 -- 10/09/2026: Sustitución electrofílica Aromática}
\subsubsection*{\textcolor{blueunam}{Clase 8 -- 10/09/2026: Sustitución electrofílica Aromática}}

\textbf{Convención de flechas curvas (eventos electrónicos):} las flechas curvas 
representan el movimiento de electrones, no de átomos. Siempre salen desde donde 
hay densidad electrónica (par libre, enlace $\pi$, enlace $\sigma$) y apuntan hacia 
donde hay deficiencia electrónica. Nunca al revés. Una flecha representa el 
movimiento de un par de electrones. \\

En reacciones entre especies, el nucleófilo (densidad negativa, par libre) 
ataca al electrófilo (deficiencia electrónica, carga positiva o $\delta^{\color{catred}+}$). 
La flecha sale del nucleófilo y apunta al electrófilo. Ejemplo: en 
\ce{OH^{\color{anblue}-} + H^{\color{catred}+}}, la flecha sale del par libre del \ce{O^{\color{anblue}-}} y apunta al \ce{H^{\color{catred}+}}.
La atracción es mutua por electrostática, pero la convención sigue 
el movimiento de los electrones, no del ``hueco'' electrónico. \\

\textbf{Sustitución electrofílica Aromática} \\

\begin{center}
\schemestart
\chemfig{*6(=-=-(-\charge{0=\:,180=\:}{O}-[:40]H)=-)}
\+
$E^{\color{catred}+}$
\arrow{}
\chemname{\chemfig{*6(=-=-(=\charge{180=\:}{\textcolor{catred}{O^{\color{catred}+}}}-[:40]H)-(-[:180]H)(-[:140]E)-)}}{Complejo sigma}
\arrow{}
\chemfig{*6(=-=-(-\charge{0=\:,180=\:}{O}-[:40]H)=(-E)-)}
\+
$H^{\color{catred}+}$
\schemestop
\end{center}

La \textbf{Sustitución Electrofílica Aromática (SEA)} consiste en sustituir 
un hidrógeno del anillo por un electrófilo entrante $E^{\color{catred}+}$, conservando la 
aromaticidad al final. El mecanismo ocurre en dos pasos: primero el sistema 
$\pi$ del anillo ataca al $E^{\color{catred}+}$ formando el complejo sigma (el anillo pierde 
aromaticidad temporalmente), luego se pierde el $H^{\color{catred}+}$ del carbono atacado 
restaurando la aromaticidad.

En el ejemplo, el \ce{OH} es un DG que dirige a \textit{orto} y \textit{para}. Se muestra la 
sustitución en \textit{orto}, pero en la práctica el producto mayoritario es el \textit{para} 
por menor impedimento estérico del grupo \ce{OH}. En reacciones más complejas 
o cuando se busca controlar el producto obtenido, conocer la regioquímica 
(\textit{orto}, \textit{meta}, \textit{para}) y el efecto del impedimento estérico es indispensable para 
predecir qué compuesto predominará. \\

\textbf{Más ejemplos de SEA}

\textbf{Nitración}
\begin{center}
\schemestart
\chemfig{*6(=-=-=-)}
\arrow{->[\ce{HNO3}][\ce{H2SO4}, $\Delta$]}[0,1.5]
\chemfig{*6(=-=(-NO_2)-=-)}
\schemestop
\end{center}

\textbf{Sulfonación}
\begin{center}
\schemestart
\chemfig{*6(=-=-=-)}
\arrow{->[\ce{SO3}][\ce{H2SO4}]}[0,1.5]
\chemfig{*6(=-=(-S(=[:120]O)(=[:300]O)(-OH))-=-)}
\schemestop
\end{center}

\textbf{Halogenación}
\begin{center}
\schemestart
\chemfig{*6(=-=(-NO_2)-=-)}
\arrow{->[1-\ce{Cl2}, \ce{AlCl3}][2-\ce{Br2}, \ce{FeBr3}]}[0,1.5]
\chemfig{*6(=-=(-Cl)-=-)}, 
\chemfig{*6(=-=(-Br)-=-)}
\schemestop
\end{center}

\textbf{Alquilación}
\begin{center}
\schemestart
\chemfig{*6(=-=-=-)}
\arrow{->[\ce{R-X}][\ce{AlCl3}]}[0,1.5]
\chemfig{*6(=-=(-R)-=-)}
\schemestop
\end{center}

\textbf{Acilación}
\begin{center}
\schemestart
\chemfig{*6(=-=-=-)}
\arrow{->[\chemfig{R-[:30](=[:90]O)-[:330]Cl}][\ce{AlCl3}]}[0,1.5]
\chemfig{*6(=-=(-[:30](=[:90]O)-[:330]R)-=-)}
\schemestop
\end{center}

\textbf{Sandmeyer}
\begin{center}
\schemestart
\chemfig{*6(=-=(-NO_2)-=-)}
\arrow{->[\ce{Sn}][\ce{HCl}]}[0,1.5]
\chemfig{*6(=-=(-NH_2)-=-)}
\arrow{->[\ce{HCl}][\ce{NaNO2}]}[0,1.5]
\schemestop
\end{center}
\begin{center}
\schemestart
\chemname[-1cm]{\chemfig{*6(=-=(-N^{\textcolor{catred}{+}}~[:0]N)-=-)}}{Sal de diazonio}
\+
$\text{Cl}^{\textcolor{anblue}{-}}$
\arrow{->[\ce{CuX}][\ce{HX}]}[0,1.5]
\chemfig{*6(=-=(-X)-=-)}
\schemestop
\end{center}
Donde \ce{X} = \ce{Cl}, \ce{Br}, \ce{CN}, \ce{I} \\

\phantomsection
\addcontentsline{toc}{subsubsection}{Clase 9 -- 24/09/2026: Clasificación de las reacciones químicas inorgánicas}
\subsubsection*{\textcolor{blueunam}{Clase 9 -- 24/09/2026: Clasificación de las reacciones químicas inorgánicas}}

\begin{itemize}
    \item \textbf{Combinación:} dos o más sustancias se unen para formar un único producto.
    \begin{center}
    \schemestart
    \chemfig{A} \+ \chemfig{B} \arrow{->} \chemfig{A-B}
    \schemestop
    \end{center}
    Ejemplo: \ce{2H2 + O2 -> 2H2O}

    \item \textbf{Descomposición:} una sustancia se rompe en dos o más productos más simples.
    \begin{center}
    \schemestart
    \chemfig{A} \arrow{->} \chemfig{B} \+ \chemfig{C}
    \schemestop
    \end{center}
    Ejemplo: \ce{2NaN3 -> 3N2 + 2Na}

    \item \textbf{Desplazamiento:} un elemento más reactivo sustituye a otro dentro de un compuesto.
    \begin{center}
    \schemestart
    \chemfig{A-B} \+ \chemfig{C} \arrow{->} \chemfig{A-C} \+ \chemfig{B}
    \schemestop
    \end{center}
    Ejemplo: \ce{Zn + CuSO4 -> ZnSO4 + Cu}
\end{itemize}

\textbf{Nota:} Las descomposiciones orgánicas son raras, pero pueden darse, como en las aminas.
\begin{center}
\schemestart
\chemfig{*6(=-=(-N=[:0]N^{\color{catred}+}=[:0]N^{\color{anblue}-})-=-)} \arrow{->} \chemfig{2N_2} \+ \chemfig{*6(=-=(-[:30]N-[:330]N-[:30]*6(=-=-=-))-=-)}
\schemestop
\end{center}

Por intercambio de partícula: \\
\begin{itemize}
    \item Ácido-base
    \begin{itemize}
        \item Brønsted --- Intercambio de protones (H$^{\color{catred}{+}}$)
        \item Lewis --- Donación de par libre (complejación)
    \end{itemize}
    \item Redox --- Intercambio de electrones (\ce{e^{\color{anblue}{-}}})
\end{itemize}

\begin{enumerate}[label=\alph*)]
    \item \schemestart
$\underbrace{\chemfig{*6(=-=-N=-)}}_{\text{Base}}$
\+
$\underbrace{\chemfig{HBr}}_{\text{Ácido}}$
\arrow{->}
$\underbrace{\chemfig{*6(=-=-N^{\color{catred}+}(-[:90]H)=-)}}_{\text{Ácido conjugado}}$
\+
$\underbrace{\chemfig{Br^{\color{anblue}-}}}_{\text{Base conjugada}}$
\schemestop : Ácido --- Base
    \item \schemestart
\chemfig{Zn^{0}}
\+
\chemfig{Cu^{\color{catred}{2+}}SO_{4}}
\arrow{->}
\chemfig{Zn^{\color{catred}{2+}}SO_{4}}
\+
\chemfig{Cu^{0}}
\schemestop : Redox
    \item \schemestart
\chemfig{NH3}
\+
\chemfig{BH3}
\arrow{->}
\chemfig{H_3N-BH_3}
\schemestop : Ácido-base de Lewis

\begin{center}
\begin{tikzpicture}
% NH3
\node (nh3) at (0,0) {%
\chemfig{\charge{90=\:}{N}(-[:210]H)(-[:330]H)(-[:90]H)}%
    };
\node[below] at (0,-1.2) {\small Base de Lewis};

% Flecha de reacción
\draw[-{Stealth}, thick] (1.5,0) -- (2.5,0);

% BH3
\node (bh3) at (4.2,0) {%
\chemfig{B(-[:30]H)(-[:150]H)(-[:270]H)}%
    };
\node[below] at (4.2,-1.2) {\small Ácido de Lewis};

% Orbital pz vacío sobre B
\orbital[pos={(4.2, 0.4)}, pcolor=gray!50, ncolor=gray!30, 
             opacity=0.8, scale=1.2]{pz}
\node[above] at (4.2, 1.4) {\small $2p_z$ vacío};

% Segunda flecha de reacción
\draw[-{Stealth}, thick] (7.2,0) -- (8.2,0);

% Producto H3N+-B-H3
\node at (11.2,0) {%
\chemfig{(N^{\color{catred}{+}}(-[:100]H)(<:[:210]H)(<[:230]H)-B^{\color{anblue}{-}}(<:[:30]H)(<[:10]H)(-[:280]H))}%
    };
\node[below] at (11.2,-1.2) {\small Complejo o aducto};

\end{tikzpicture}
\end{center}

    \item \schemestart
\chemfig{CuCl2}
\+
2\,\chemfig{NH3}
\arrow{->}
\ce{[Cu(NH3)2Cl2]}
\schemestop :\\ Es una complejación. Expliquemos cómo sucede para este caso, que se puede extrapolar a cualquier complejación \\

\vspace{1cm}

\schemestart
\chemfig{@{cl1}Cl^{\color{anblue}-}-Cu^{\color{catred}{2+}}-@{cl2}Cl^{\color{anblue}-}}
\+
\chemname{%
\begin{tikzpicture}[baseline=-0.5ex]
    \node (nh3) at (0,0) {%
        \chemfig{N(-[:200]H)(<[:250]H)(<:[:330]H)}%
    };
    \orbital[pos={(0.1, 0.3)}, color=anblue, nelec=2, opacity=0.6, scale=0.8, rotate=90]{lobe}
    \node[above] at (0.1, 1) {\small HOMO ($n$)};
\end{tikzpicture}%
}{$L^0$}
\arrow{->}
\chemfig{@{cu2}Cu^{\color{catred}{2+}}(-[:0]Cl)(-[:180]Cl)(-[:90]@{n1}NH_3)(-[:270]@{n2}NH_3)}
\schemestop
\chemmove{
    \draw[-{Stealth}, thick] (cl1) .. controls +(90:0.8cm) .. +(-0.8, 1.2)
        node[above] {\small $L^{\color{anblue}-}$};
    \draw[-{Stealth}, thick] (cl2) .. controls +(90:0.8cm) .. +(0.8, 1.2)
        node[above] {\small $L^{\color{anblue}-}$};
    \draw[-{Stealth}, black] (cu2) -- +(2.0, 0.6)
        node[right]{\small Centro metálico};
    \draw[-{Stealth}, black] ($0.5*(cu2)+0.5*(n1)$) -- +(1, 0.8)
        node[right]{\small Enlace de coordinación};
    \draw[-{Stealth}, black] ($0.5*(cu2)+0.5*(n2)$) -- +(1,-0.8)
        node[right]{\small Enlace de coordinación};
} \\

El \ce{Cu^{2+}} tiene orbitales $d$ con espacio disponible para aceptar 
pares de electrones. El \ce{NH3} dona su par libre formando un enlace 
dativo: ambos electrones provienen del ligante. Como no hay transferencia 
real de electrones, el estado de oxidación del Cu no cambia. \\

\textbf{Nota:} Todos los metales de transición son del bloque $d$, pero 
no todos los del bloque $d$ son metales de transición. El Zn, Cd y Hg 
tienen la subcapa $d$ completamente llena (d$^{10}$) en todos sus estados 
de oxidación comunes, por lo que quedan excluidos por definición IUPAC.


    \item \schemestart
\ce{PdCl2}
\+
\ce{2PPh3}
\arrow{->}
\ce{PdCl2(PPh3)2}
\schemestop :\\ Es una complejación. Expliquemos cómo sucede: \\
\textbf{Nota:} El Ph es un fenilo igual se puede escribir como $\varnothing$ :
\schemestart
\chemfig{*6(=-=(-)-=-)}
\schemestop
\vspace{1cm}

\schemestart
\chemfig{@{cl1}Cl^{\color{anblue}-}-Pd^{\color{catred}{2+}}-@{cl2}Cl^{\color{anblue}-}}
\+
\chemname{%
\begin{tikzpicture}[baseline=-0.5ex]
    \node (nh3) at (0,0) {%
        \chemfig{P(-[:200]Ph)(<[:250]Ph)(<:[:330]Ph)}%
    };
    \orbital[pos={(0.1, 0.3)}, color=anblue, nelec=2, opacity=0.6, scale=0.8, rotate=90]{lobe}
    \node[above] at (0.1, 1) {\small HOMO ($n$)};
\end{tikzpicture}%
}{$L^0$}
\arrow{->}
\chemfig{@{cu2}Pd^{\color{catred}{2+}}(-[:0]Cl)(-[:180]Cl)(-[:90]@{n1}PPh_3)(-[:270]@{n2}PPh_3)}
\schemestop
\chemmove{
    \draw[-{Stealth}, thick] (cl1) .. controls +(90:0.8cm) .. +(-0.8, 1.2)
        node[above] {\small $L^{\color{anblue}-}$};
    \draw[-{Stealth}, thick] (cl2) .. controls +(90:0.8cm) .. +(0.8, 1.2)
        node[above] {\small $L^{\color{anblue}-}$};
    \draw[-{Stealth}, black] (cu2) -- +(2.0, 0.6)
        node[right]{\small Centro metálico};
    \draw[-{Stealth}, black] ($0.5*(cu2)+0.5*(n1)$) -- +(1, 0.8)
        node[right]{\small Enlace de coordinación};
    \draw[-{Stealth}, black] ($0.5*(cu2)+0.5*(n2)$) -- +(1,-0.8)
        node[right]{\small Enlace de coordinación};
} \\
\end{enumerate}

\textbf{Nota:} Todo lo que rodea al metal es la esfera de coordinación \\
\textbf{Nota sobre el par libre del N y P:} La notación $2p_z$ sería 
una simplificación incorrecta. El N en \ce{NH3} tiene configuración 
$1s^2 2s^2 2p^3$ y forma 3 enlaces, por lo que el par libre proviene 
de electrones originalmente en el $2s$, no en un $2p$ puro. En el modelo 
de hibridación sp$^3$ (modelo matemático, no entidad física real) el par 
libre ocupa un orbital híbrido. En TOM, lo correcto es identificarlo como 
el \textbf{HOMO} (\textit{Highest Occupied Molecular Orbital}, el orbital 
lleno de mayor energía; véase \hyperref[sec:anexo:tom:frontera]{Anexo A.1.12}), 
orbital no enlazante ($n$): tiene predominantemente carácter $s$ con 
contribución $p$, forma asimétrica similar a un lóbulo $p$ pero más 
compacta. Se representa como un lóbulo tipo $p$ por convención didáctica, 
no porque sea un $p$ puro. Lo mismo aplica para el P en \ce{PPh3} con su 
orbital $3s$ de origen, cuyo HOMO ($n$) tiene aún más carácter $s$ que el 
del N, lo que lo hace más esférico y menos disponible para resonancia con 
sistemas $\pi$, pero mejor ligante por su mayor polarizabilidad.